\documentclass[12pt,a4paper]{article}
\usepackage[utf8]{inputenc}
\usepackage[T1]{fontenc}
\usepackage{amsmath}
\usepackage{amsthm}
\usepackage{amsfonts}
\usepackage{amssymb}
\usepackage{lmodern}
\usepackage{graphicx}
\usepackage{booktabs}
\usepackage[left=2cm,right=2cm,top=2cm,bottom=2cm]{geometry}
\author{Filippo Perugini}
\title{DRAFT! Exact Calibration of\\ Multi-Factor Stochastic Interest Rates \\
Local-Stochastic Volatility  (SIR-LSV) models \\ with Particle Algorithm and extensions to Long Dated FX}
\newcommand\numberthis{\addtocounter{equation}{1}\tag{\theequation}}

\theoremstyle{plain}
\newtheorem{thm}{Theorem}[section]
\newtheorem{lem}[thm]{Lemma}
\newtheorem{prop}[thm]{Proposition}
\newtheorem*{cor}{Corollary}

\theoremstyle{definition}
\newtheorem{defn}{Definition}[section]
\newtheorem{conj}{Conjecture}[section]
\newtheorem{exmp}{Example}[section]

\theoremstyle{remark}
\newtheorem*{rem}{Remark}
\newtheorem*{sol}{Solution}
\newtheorem*{exe}{Exercise}

\newtheorem*{note}{Note}

\makeatletter
\def\th@plain{%
  \thm@notefont{}% same as heading font
  \itshape % body font
}
\def\th@definition{%
  \thm@notefont{}% same as heading font
  \normalfont % body font
}
\makeatother

\newcommand{\MD}{\text{\DH}}
\date{Draft of 2016, revised November 2026}

\begin{document}

\maketitle

\begin{abstract}
We derive the particle (McKean) calibration of local volatility models with stochastic interest rates, following Guyon and Henry-Labord\`ere, and extend it to long-dated FX with Hull-White domestic and foreign rates. The three-factor calibration condition is written under a single measure, so that one system of particles suffices. Its Malliavin representation is obtained without solving the coupled two-parameter system of Malliavin derivatives: integrating the derivatives against the Hull-White kernels yields four auxiliary processes per particle, simulated forward together with the state, and a correction to the Dupire local variance which involves neither $\partial_K^2C$ nor tail expectations. The representation reduces to the known one-factor formula, matches the Gaussian closed form, and its building blocks are verified by Monte Carlo. In a particle calibration to CME EUR/USD option settlements with illustrative rate parameters, the repricing error attributable to the stochastic rates is at most 0.005 volatility points over the quoted expiries, against 0.19 and 0.47 for the two parameter sets when the correction is omitted. For a prescribed volatility, the discounted density and the first moments of the auxiliary processes solve a closed system of five linear forward equations, which gives a deterministic formulation of the calibration. We give a conditional local existence result for flat-smile inputs, an input-compatibility obstruction, and a separate well-posedness theorem for a regularized system; existence for general smile inputs remains open in this work.
\end{abstract}

\tableofcontents

\section{Introduction}\label{sec: introduction}
While for short-dated options (less than 1 year), assuming constant interest rates does not lead to
significant mispricing, for long-dated options the effect of interest rate volatility becomes increasingly
pronounced with increasing maturity and can become as important as that of the FX spot volatility.

Here we consider a three factor model in which uncertainty in the spot process is explained by a geometric Brownian motion and both the domestic and foreign interest rates follow Hull-White processes.
To incorporate smile/skew effects in the three-factor pricing model, we introduce a Local Volatility function in the spot process.

The goal of this section is to explicitly derive the Particle Algorithm of \cite{guyonLabordere} in a Three Factor Local Volatility Stochastic Interest Rate (SIR-LV) framework. We derive the calibration condition for one factor stochastic domestic and foreign rates with unspecified dynamics, and then, for Hull-White rates, a Malliavin representation of it which can be evaluated with a single system of particles.

Monte Carlo calibration of local volatility in an FX setting with two stochastic rates is studied in \cite{ogetbilGanesanHientzsch} (see also \cite{ogetbilHientzsch} for the corresponding Dupire formula), and a particle method with control variates for a four-factor FX local-stochastic volatility model in \cite{cozma}; an optimal transport approach to stochastic-rates calibration is developed in \cite{josephLoeperObloj}.

\section{Malliavin calculus: the tools used}\label{sec: tools}
Before moving to the actual computation of the Malliavin derivatives in (\ref{eq: threeFactorGeneralMalliavin}) required for an explicit formulation of the Particle Algorithm in the case of three-factor FX SIR-LV model, I attempt to give a short introduction to the tools needed to understand the following derivations.
I need to assume familiarity with the basics of Malliavin calculus. These can be found in the first chapter of \cite{nualart} and I will refer to the Nualart's book through out the computations.

The results I am going to use are the \textit{commutative relation} between the Derivative and Divergence operators (Proposition 1.3.8 in \cite{nualart}) and the \textit{chain rule} for the Derivative operator (Proposition 1.2.3 in \cite{nualart}).

\begin{prop}[Commutative Relation between $\MD$ and $\delta$]\ \\
Provided that the condition of Proposition 1.3.8 in \cite{nualart} are satisfied, one has:

\begin{equation}\label{eq: commumativeRelation}
\MD_s( \delta(u) )= u_s\mathbb{I}_{s\leq t} + \int_0^t\MD_s u_r \mathbb{I}_{s\leq t}dW_r
\end{equation}
for a stochastic process $u$ on $L^2([0,t] \times \Omega)$.
\end{prop}

To get familiar with such tools, I am going to solve the following easy exercise (exercise 2.2.1 in\cite{nualart} ):\\

\begin{exe}
Let $\sigma$ and $b$ be continuously differentiable functions on $\mathbb{R}$ with bounded
derivatives. Consider the solution $X = \{X_t, t \in [0, T]\}$ of the stochastic
differential equation
$$ X_t = x_0 + \int_0^t\sigma(X_r)dW_r + \int_0^tb(X_r)dr$$
Show that for $s\leq t$ we have
$$ \MD_sX_t = \sigma(X_s)\exp\bigg(\int_s^t\sigma'(X_r)dW_r  + \int_s^t[b'(X_r)-\frac{1}{2}\sigma'(X_r)^2]dr  \bigg) $$
Where $\MD$ is the Malliavin derivative with respect to the Brownian Motion $W$.
\end{exe}\
\\
\begin{sol}
Let's first rewrite the previous SDE as:
$$ X_t = x_0 + \delta\big(\sigma(X)\big) + <b(X),\mathbb{I}_{[0,t]}> $$
where $\delta$ is the Skorohod integral and $<\cdot, \cdot>$ is the inner product, both under the White Noise case (Section 1.2.1 in \cite{nualart}).\\
Then:
\begin{align*}
\MD_sX_t = & \MD_s\delta\big(\sigma(X)\big) + <\MD_sb(X),\mathbb{I}_{[0,t]}>\\
= & \sigma(X_s)\mathbb{I}_{s\leq t} + \delta\big(\MD_s\sigma(X)\mathbb{I}_{s\leq t}\big) +  <\MD_sb(X),\mathbb{I}_{[0,t]}>
\end{align*} 
Where we get the second line from an application of the \textit{commutative relation}.
By applying the \textit{chain rule} of the Malliavin derivative to $\sigma$ and $b$ we get:
\begin{align*}
\MD_sX_t =\MD_sX_t\mathbb{I}_{s\leq t}= & \sigma(X_s)\mathbb{I}_{s\leq t} + \delta\big( \sigma'(X)\cdot \MD_sX \mathbb{I}_{s\leq t}\big) + < b'(X) \cdot\MD_s X,\mathbb{I}_{[0,t]}>\\
= & \sigma(X_s) + \int_0^t \sigma'(X_r)\cdot\MD_sX_r \mathbb{I}_{s\leq t} dW_r + \int_0^t b'(X_r) \cdot\MD_s X_r\mathbb{I}_{s\leq t} dr
\end{align*} 
Define now $S_t:= \MD_sX_t\mathbb{I}_{s\leq t}$, one may write
$$ S_t = \sigma(X_s)\mathbb{I}_{s\leq t} + \int_0^t \sigma'(X_r)S_rdW_r + \int_0^t b'(X_r) S_rdr,$$
thus
$$ dS_t = \sigma'(X_t)S_tdW_t + b'(X_t) S_t dt$$
which is the SDE driving a Geometric Brownian Motion with initial value equal to $\sigma(X_s)$. So the result is proven $\diamond$ .
\end{sol}

The exercise is the scalar case of a general fact used below: for a diffusion $X$, $\MD_sX_t=J_tJ_s^{-1}\times$(diffusion coefficient at $s$), where $J$ is the first variation (tangent) process of $X$. For a system, $J$ is matrix-valued.

\section{Background: one stochastic rate}\label{sec: oneRate}
One of the great features of the particle method is that it easily extends to
the case where interest rates are also stochastic. An LSVM with stochastic
interest rates (SIR-LSVM) is defined in a risk-neutral measure  $ \mathbb{Q}  $ by:

\begin{equation}  \label{eq:spot} dS_t=r_tS_t dt + \sigma(t,S_t)a_t S_tdW_t 
\end{equation}

where the short term rate $r_t$ and the stochastic volatility $a_t$ are Ito processes.
\subsection{The Calibration Condition}
In this case the calibration condition reads:

\begin{equation} \label{eq:calibration_condition_spot} \sigma(t,K)^2 \frac{\mathbb{E}^{\mathbb{Q}}[ D_{0t}a_t^2 | S_t = K ]}{ \mathbb{E}^{\mathbb{Q}}[ D_{0t} | S_t = K ] } = \sigma_{Dup}(t,K) ^2 - \frac{ \mathbb{E}^{\mathbb{Q}}[ D_{0t}(r_t-r^0_t)\mathbb{I}_{S_t>K} ]}{ \frac{1}{2}K\partial^2_KC(t,K) },
\end{equation}

where the discount factor $ D_{0t}$ is given by $ D_{0t}=exp\bigg(-\int_0^t r_sds\bigg)$ and

$r_t^0 = \mathbb{E}^{\mathbb{Q}^t} [r_t] = -\partial_t\log P_{0t}$, with $ P_{st} $ the time-$s$ value of the bond of maturity $t$.\\

From the calibration condition (\ref{eq:calibration_condition_spot}), the dynamics of the calibrated SIR-LSVM read as the following nonlinear McKean diffusion: 

$$ dS_t=r_tS_t dt + \sigma(t,S_t,\mathbb{Q}_t)a_t S_tdW_t .$$

Alternatively, it reads as the following nonlinear McKean diffusion for the
forward $f_t = S_t/P_{tT}$ in the forward measure $\mathbb{Q}^T$ , where $T$ denotes the last maturity date for which we want to calibrate the market smile:

\begin{equation}\label{eq: forwardPrice}
\frac{df_t}{f_t}= \sigma(t,P_{tT}f_t,\mathbb{Q}_t^T)a_t dW_t^T -\sigma_P^T(t)dB^T_t
\end{equation} 

where 

\begin{equation}\label{eq: localVolatility_forward}
\sigma(t,K,\mathbb{Q}_t^T)^2=\bigg( \sigma_{Dup}(t,K)^2-P_{0T} \frac{ \mathbb{E}^{\mathbb{Q}^T}[ P_{tT}^{-1}(r_t-r^0_t)\mathbb{I}_{S_t>K} ]}{ \frac{1}{2}K\partial^2_KC(t,K) } \bigg)\cdot  \frac{\mathbb{E}^{\mathbb{Q}^T}[P_{tT}^{-1} | S_t = K ]}{ \mathbb{E}^{\mathbb{Q}^T}[ P_{tT}^{-1}a_t^2 | S_t = K ] }
\end{equation}

$ \sigma_P^T(t) $ is the volatility of the bond $ P_{tT}$, and $B^T_t$ the (possibly multi-dimensional) $\mathbb{Q}^T$-Brownian motion that drives the interest rate curve.

\subsection{The Particle Method}
A \textit{particle} is then described by the vector-valued process $(f_t^i,a^i_t,r_t^i)$, for $1\leq i \leq N$, where $N$ is the number of simulated particles.
Using the representation (\ref{eq: localVolatility_forward}) for the local volatility, we then define:

\begin{equation}\label{eq: discrete_localVol_general}
\sigma_N(t,K)^2= \bigg( \sigma_{Dup}(t,K)^2 - P_{0T}\frac{\frac{1}{N}\sum_{i=1}^N (P_{tT}^{i})^{-1}(r_t^i-r_t^0)\mathbb{I}_{S_t^i>K} }{\frac{1}{2}K\partial^2_KC(t,K)}\bigg)\cdot\frac{\sum_{i=1}^N (P_{tT}^{i})^{-1} \delta_{t,N}(S_t^i-K)}{\sum_{i=1}^N (P_{tT}^{i})^{-1} (a_t^i)^2\delta_{t,N}(S_t^i-K)}
\end{equation}

and simulate:
\begin{equation}\label{eq: discrete_forwardPrice}
df_t^i = f_t^i\sigma_N(t,f_t^i P_{tT}^i)a_t^idW_t^i-f_t^i\sigma_P^{T,i}(t)dB_t^i
\end{equation}

\subsection{Acceleration Techniques}\label{subsec: acceleration}
\begin{itemize}
\item Not necessary to compute the sum in (\ref{eq: discrete_localVol_general}) for all $i$
\item Save time: compute $\sigma_N$ for a grid $G_{K,t}$ of values of $K$ with size $N_{K,t}<N$
\item Then inter- and extra-polate (cubic splines)
\end{itemize}

\smallskip
\begin{itemize}
\item Moreover: a large number of terms in the sum in (\ref{eq: discrete_localVol_general}) is negligible
\item Disregard $f_t^i$. when it is far from $K$, i.e., $\delta_{t,N}(S_t^i-K) < \eta$ ($\eta$ being a threshold)
\item This requires sorting the particles according to the spot value
\item The cost $O(N\ln N)$ is compensated by the acceleration with the grid
\end{itemize}

\subsection{Malliavin representation of the local volatility}\label{subsec: malliavinGeneral}
Thanks to an application of the Clark-Ocone's formula we can give another expression of the contribution of stochastic interest rates to local volatility:
\begin{equation}\label{eq: malliavinRepresentation}
\frac{ \mathbb{E}^{\mathbb{Q}}[ D_{0t}(r_t-r^0_t)\mathbb{I}_{S_t>K} ]}{ \frac{1}{2}K\partial^2_KC(t,K) } = \frac{2}{K}\int _0^t \mathbb{E}^{\mathbb{Q}^t}[\sigma_r^t(s)\cdot \MD _s^{B^t}S_t  |S_t = K ]ds
\end{equation}

where $ \MD_s^{B^t} $ is the Malliavin derivative with respect to the Brownian motion $B^t$.
This brings few different benefits:
\begin{itemize}
\item[1.] instead of an unconditional expectation we get a conditional expectation given $ S_t = K $ (easier for MC);
\item[2.] the term $\partial_K^2C$ cancels out: fortunate case as it is sensitive to the strike interpolation/extrapolation method;
\item[3.] both $\mathbb{E}^{\mathbb{Q}}[ D_{0t}(r_t-r^0_t)\mathbb{I}_{S_t>K} ]$ and $K\partial_K^2C$ are very small for strikes $K$ that are far from the money. Hence we get rid of a problematic $0/0$ ratio.
\end{itemize}
By specifying a particular model for the interest rate and computing the Malliavin derivative in (\ref{eq: malliavinRepresentation}), one gets more explicit formulas for the interest rate contribution to the local volatility.
We shall see this applied to case of the Hull-White model.

\subsection{The hybrid Hull-White/Dupire model}\label{subsec: hwDupire}
Assume now that $a_t \equiv 1$. We consider the particular case where $r_t$ follows the Hull-White model:

\begin{equation}\label{eq: hull_white}
dr_t = \kappa\big(\theta(t)-r_t\big)dt+\sigma_r dB_t
\end{equation}
with $B$ is $\mathbb{Q}$-Brownian motion, correlated with $W$, the correlation being $\rho$.

The function $ \theta(t)$ is chosen such that the zero coupon curve observed at $t = 0$ is matched:
$$\theta(t) = \frac{1}{\kappa}\frac{\partial f^{mkt}}{\partial t}(0,t)+f^{mkt}(0,t)+\frac{\sigma_r^2}{2\kappa^2}(1-e^{-2\kappa t})  $$
where
$$f^{mkt}(0,t)=-\frac{\partial\ln P_{0t}^{mkt}}{\partial t} = r_t^0$$

The bond prices are given by:
\begin{equation}\label{eq: bondPrice_hullWhite}
P_{tT}= \frac{P_{0T}^{mkt}}{P_{0t}^{mkt}}\exp\bigg( -A(t,T)(r_t-r_t^0)-\frac{\sigma_r^2}{4\kappa^3}(e^{-\kappa T}-e^{-\kappa t})^2(e^{2\kappa t}-1)     \bigg)
\end{equation}
where
$$A(t,T)= \frac{1-e^{ -\kappa(T-t) }}{\kappa}$$
and the volatility of the zero-coupon bond is 
$$\sigma_P^T(t)=-\sigma_r A(t,T)$$\\

In view of simulating all processes under the terminal forward measure $\mathbb{Q}^T$ ,
we aim at writing $r_t -r_t^0$ in terms of the $\mathbb{Q}^T$-Brownian motion
$$ B^T_t=B_t-\int_0^t\sigma_P^T(s)ds=B_t+\sigma_r\int_0^tA(s,T)ds $$

After computations the explicit solution to (\ref{eq: hull_white}) under the $\mathbb{Q}^T$-measure is:
\begin{equation}\label{eq: hullWhiteSolution}
r_t -r_t^0= e^{-\kappa t}\bigg(\sigma_r\int_0^t e^{\kappa s}dB_s^T-\frac{\sigma_r^2}{\kappa^2}\sinh(\kappa t)\big(1-e^{-\kappa(T-t)} \big)   \bigg)
\end{equation}

\subsubsection{Malliavin representation under Hull-White}\label{subsec: malliavin1F}
Under these assumptions, (\ref{eq: malliavinRepresentation}) reads:

\begin{equation}\label{eq: malliavinHullWhite}
\frac{ \mathbb{E}^{\mathbb{Q}}[ D_{0t}(r_t-r^0_t)\mathbb{I}_{S_t>K} ]}{ \frac{1}{2}K\partial^2_KC(t,K) } = 
2\sigma_re^{-\kappa t}\mathbb{E}^{\mathbb{Q}^t}[ V_t(\rho U_t + \Theta_t\Xi_t-\Lambda_t)|S_t = K]
\end{equation}

where:
\begin{align*}
  \frac{dV_t}{V_t} &= S_t\partial_S\sigma(t,S_t)\{dW_t^T + [\rho\sigma_P^T(t,r_t)- \sigma(t,S_t) ]dt\} ,\ \ V_0=1\\
  U_t &= \int_0^te^{\kappa s}\frac{\sigma(s,S_s)}{V_s}ds\\ \numberthis \label{eq: malliavinWeights}
  \Theta_t &= \int_0^t\frac{e^{-\kappa s}}{V_s}ds \\  
  \Lambda_t &= \int_0^t\sigma_re^{2\kappa s}\Theta_s ds\\
  \Xi_t &= \sigma_r\frac{e^{2\kappa t}-1}{2\kappa}
\end{align*}

The computation of (\ref{eq: malliavinHullWhite}) for all $t$, requires the simulation of only 3 processes
$f_t$, $r_t$, $V_t$, and 3 integrals $U_t$, $\Theta_t$, $\Lambda_t$.
The local volatility (\ref{eq: localVolatility_forward}) can be now rewritten as:

\begin{align*}
\sigma(t,K,\mathbb{Q}_t^T)^2=  \sigma_{Dup}(t,K)^2 &-2\sigma_r\rho e^{-\kappa t}\frac{\mathbb{E}^{\mathbb{Q}^T}[P_{tT}^{-1}V_tU_t|S_t=K]}{\mathbb{E}^{\mathbb{Q}^T}[P_{tT}^{-1}|S_t=K]} \\ 
 &  -2\sigma_r e^{-\kappa t}\frac{\mathbb{E}^{\mathbb{Q}^T}[P_{tT}^{-1}V_t(\Theta_t\Xi_t-\Lambda_t)|S_t=K]}{\mathbb{E}^{\mathbb{Q}^T}[P_{tT}^{-1}|S_t=K]} \numberthis \label{eq: localVolatility_forward_HullWhite}
\end{align*}

\subsubsection{The particle method}
In the particle algorithm for the hybrid Hull-White/Dupire model, a particle can then be described by $(f_t, r_t- r^0_t )$. Simulate the $T$-forward price $f_t = S_t/P_{tT}$ under the final forward measure as follows:
$$ df_t^i=  f_t^i\sigma_N(t,f_t^iP_{tT}^i)dW_t^i+f_t^i\sigma_r A(t,T)dB_t^i $$
where $B^i$ and $W^i$ are $\mathbb{Q}^T$-Brownian motions.\\
Using the representation (\ref{eq: localVolatility_forward_HullWhite}) for the local volatility, we define:

\begin{align*}
\sigma_N(t,K)^2 = \sigma_{Dup}(t,K)^2 & -2\sigma_r\rho e^{-\kappa t}\frac{\sum_{i=1}^N \Pi_t^{i} V_t^iU_t^i \delta_{t,N}(S_t^i-K) }{\sum_{i=1}^N \Pi_t^{i} \delta_{t,N}(S_t^i-K)} \\
&  -2\sigma_r e^{-\kappa t}\frac{\sum_{i=1}^N \Pi_t^{i}V_t^i(\Theta^i_t\Xi^i_t-\Lambda^i_t)\delta_{t,N}(S_t^i-K) }{\sum_{i=1}^N \Pi_t^{i} \delta_{t,N}(S_t^i-K)} \numberthis \label{eq: discrete_localVolatility_forward_HullWhite}
\end{align*}

with:

$$ \Pi_t^{i} = \exp\bigg( A(t,T)(r_t^i-r_t^0)+\frac{\sigma_r^2}{4\kappa^3}(e^{-\kappa T}-e^{-\kappa t})^2(e^{2\kappa t}-1)     \bigg) $$

The time discretisation of these processes and the resulting algorithm are given in Appendix \ref{app: oneRate}.

\subsection{An equivalent one-process form}\label{subsec: oneProcess}
The five objects $V_t,U_t,\Theta_t,\Lambda_t,\Xi_t$ can be replaced by a single process. Define
$$ \ell_t := \sigma_r e^{-\kappa t}\,V_t\big(\rho U_t+\Theta_t\Xi_t-\Lambda_t\big). $$
Since $\rho U_t+\Theta_t\Xi_t-\Lambda_t$ has finite variation, with
$d(\rho U_t+\Theta_t\Xi_t-\Lambda_t)=V_t^{-1}\big(\rho e^{\kappa t}\sigma(t,S_t)+e^{-\kappa t}\Xi_t\big)dt$,
It\^o's formula gives
\begin{equation}\label{eq: oneProcess}
d\ell_t=\ell_t\Big[g(t,S_t)\big(dW_t-\sigma(t,S_t)dt\big)-\kappa\,dt\Big]
+\Big[v(t)+\rho\,\sigma_r\,\sigma(t,S_t)\Big]dt,\qquad \ell_0=0,
\end{equation}
where
$$ g(t,S)=S\,\partial_S\sigma(t,S),\qquad v(t)=\sigma_r^2\,\frac{1-e^{-2\kappa t}}{2\kappa}, $$
and $W$ is the $\mathbb{Q}$-Brownian motion, $dW_t=dW^T_t+\rho\,\sigma_P^T(t)dt$. Note that $v(t)$ is the variance of the Hull-White state at time $t$.
The right hand side of (\ref{eq: malliavinHullWhite}) is then $2\,\mathbb{E}^{\mathbb{Q}^t}[\ell_t\,|\,S_t=K]$ and (\ref{eq: discrete_localVolatility_forward_HullWhite}) becomes
\begin{equation}\label{eq: oneProcessParticle}
\sigma_N(t,K)^2=\sigma_{Dup}(t,K)^2-2\,\frac{\sum_{i=1}^N\Pi^i_t\,\ell^i_t\,\delta_{t,N}(S^i_t-K)}{\sum_{i=1}^N\Pi^i_t\,\delta_{t,N}(S^i_t-K)} .
\end{equation}
One Euler step of (\ref{eq: oneProcess}) per particle replaces the simulation of $V_t$ and the three quadratures of the previous section. This is the form that extends to the three-factor model of Section \ref{sec: threeFactor}: there, $\ell_t$ is the logarithmic sensitivity of the spot to a perturbation of the rate driver along the Hull-White kernel.

\section{The three-factor FX model}\label{sec: threeFactor}
\subsection{The Three Factor model with Local Volatility }
The model is fully described by the following SDEs:
\begin{align*}
 dS_t   = &\ (r^d_t - r^f_t)S_t dt + \sigma(t,S_t)S_tdW^S_t \\
 dr_t^d  = &\  \mu^d_t dt + \sigma_t^d dB_t^d \numberthis \label{eq: threeFactorModel}\\
 dr_t^f  = & \ \big(\mu^f_t - \rho_{Sf}\sigma(t,S_t)\sigma_t^f \big)dt + \sigma_t^f dB_t^f
\end{align*}
where $S_t$ is the spot, $r_t^d$ and $r_t^f$ the domestic and foreign interest rates respectively.\\
The volatility of the spot FX rate is the local volatility function $\sigma(t,S_t)$.
At this stage the dynamics for the two interest rates are left unspecified to let the framework as general as possible. The stochastic vector $(W^S_t, B_t^d, B_t^f)$ is a Brownian Motion under the domestic risk-neutral measure $\mathbb{Q}^d$ with correlation matrix
\[R=\left( \begin{array}{ccc}
1 & \rho_{Sd} & \rho_{Sf} \\
\rho_{Sd} & 1 & \rho_{df} \\
\rho_{Sf} & \rho_{df} & 1 \end{array} \right)\]
The term $-\rho_{Sf}\sigma(t,S_t)\sigma^f_t$ in the drift of the foreign rate is the quanto adjustment: under the foreign risk-neutral measure $\mathbb{Q}^f$ the foreign rate has drift $\mu^f_t$.

\paragraph{Measures.} The domestic and foreign discount factors are $ D_{0t}^i=\exp\big(-\int_0^t r_s^ids\big)$ for $i\in\{d,f\}$. The foreign risk-neutral measure is defined by (section 2.8 and p.256 in \cite{clark})
\begin{equation}\label{eq: RNforeign}
\left.\frac{d\mathbb{Q}^f}{d\mathbb{Q}^d}\right|_{\mathcal{F}_t}=\frac{S_tD^d_{0t}}{S_0D^f_{0t}}
=\exp\bigg( \int_0^t\sigma(s,S_s) dW^S_s - \frac{1}{2}\int_0^t\sigma(s,S_s)^2 ds\bigg),
\end{equation}
so that, denoting with a tilde the Brownian motions under $\mathbb{Q}^f$,
\begin{align*}
dW^S_t = &\ d\widetilde{W}^{S}_t + \sigma(t,S_t)dt \\
dB^d_t = &\ d\widetilde{B}^{d}_t + \rho_{Sd}\sigma(t,S_t)dt  \numberthis \label{measureChange: domesticToForeign}\\
dB^f_t = &\ d\widetilde{B}^{f}_t + \rho_{Sf}\sigma(t,S_t)dt
\end{align*}
Note that the spot driver is shifted by the full $\sigma(t,S_t)$, and that the Girsanov kernel is a function of the state. We write $\mathbb{Q}^{d^T}$ and $\mathbb{Q}^{f^T}$ for the domestic and foreign $T$-forward measures, with numeraires $P^d_{tT}$ and (in foreign currency) $P^f_{tT}$.

\subsection{The Calibration Condition}
In the proposition below, we give a necessary and sufficient condition for a three factor SIR-LVm to be calibrated to a given market smile. $C(t,K)$ denotes the fair market value of a European call option with strike $K$ and maturity $t$.
From \cite{deelstraRayee}, the Dupire volatility is of the form:
\begin{equation}\label{eq: threeFactorDupireVolatility}
\sigma_{Dup}^2(t,K) =
\frac{  \partial_t C(t,K) + (r_d^0 - r_f^0)K\partial_K C(t,K)+ r_f^0 C(t,K)}
{ \frac{1}{2}K^2\partial^2_K C(t,K) }
\end{equation}
where $r_i^0 = r_i^0(t)= \mathbb{E}^{i^t} [r^i_t] = -\partial_t\ln P^{i}_{0t}$ is the market instantaneous forward rate, $\mathbb{E}^{i^t}$ being the expectation under the $t$-forward domestic and foreign measures for $i = d, i=f$ respectively.

\begin{prop}[Calibration of 3-Factor SIR-LVm to market smiles]\label{prop: calibration}\ \\
Assume the model exists with a $C^{1,2}$ call surface and sufficient
integrability for the discounted Tanaka formula below. Assume a smooth market
surface with positive strike density, matching initial and boundary conditions,
and uniqueness for the linear forward equation with the model's fixed local
variance coefficient. Under these hypotheses the model
(\ref{eq: threeFactorModel}) is exactly calibrated to the market smile if and only if
\begin{align*}
\sigma(t,K)^2 = \sigma_{Dup}(t,K)^2 & - \frac{ \mathbb{E}^d[ D_{0t}^d(r_t^d-r^0_d)\mathbb{I}_{S_t>K} ]}{ \frac{1}{2}K\partial^2_KC(t,K) }\\
 &+ S_0\frac{ \mathbb{E}^f[ D_{0t}^f(r_t^f-r^0_f)\mathbb{I}_{S_t>K} ]}{ \frac{1}{2}K^2\partial^2_KC(t,K) }\numberthis \label{eq: threeFactorCalibrationCondition}
\end{align*}
or, equivalently, under the domestic measure only,
\begin{equation}\label{eq: singleMeasure}
\sigma(t,K)^2 = \sigma_{Dup}(t,K)^2 - \frac{ \mathbb{E}^d\big[ D_{0t}^d\big(K(r_t^d-r^0_d)-S_t(r_t^f-r^0_f)\big)\mathbb{I}_{S_t>K} \big]}{ \frac{1}{2}K^2\partial^2_KC(t,K) }
\end{equation}
\end{prop}

\begin{proof}\ \\
We follow the proof of Proposition 11.1 in \cite{guyonLabordere} and show that the result is an easy generalisation of (\ref{eq:calibration_condition_spot}). \\
Define $\mathcal{P}_t\equiv D^d_{0t}(S_t-K)^+$ to be the discounted vanilla call payoff with maturity $t$ and strike $K$ under the domestic measure. By applying Ito-Tanaka's formula on such payoff, we get:

\begin{align*}
d\mathcal{P}_t  = & - r_t^d D^d_{0t}(S_t-K)^+dt + D^d_{0t}S_t\mathbb{I}_{S_t>K} \{ (r_t^d-r_t^f)dt + \sigma(t,S_t)dW^S_t \}\\
& + \frac{1}{2} D^d_{0t}\delta_K(S_t)S^2_t\sigma(t,S_t)^2dt \\
= &  - r_t^d D^d_{0t}S_t\mathbb{I}_{S_t>K}dt + r_t^d D^d_{0t}K\mathbb{I}_{S_t>K}dt + r_t^d D^d_{0t}S_t\mathbb{I}_{S_t>K}dt\\
&  -r_t^f D^d_{0t}S_t\mathbb{I}_{S_t>K}dt + D^d_{0t}\mathbb{I}_{S_t>K}\sigma(t,S_t)S_tdW^S_t+ \frac{1}{2} D^d_{0t}\delta_K(S_t)K^2\sigma(t,K)^2dt\\
= & D^d_{0t}\mathbb{I}_{S_t>K}(r_t^d K- r_t^f S_t)dt +  D^d_{0t}\mathbb{I}_{S_t>K}\sigma(t,S_t)S_tdW^S_t 
+ \frac{1}{2} D^d_{0t}\delta_K(S_t)\sigma(t,K)^2K^2 dt
\end{align*}
By taking the expectation $\mathbb{E}^d[\cdot]$ on both sides of the above equation and by assuming that $ M_t =\int_0^tD^d_{0s}\mathbb{I}_{S_s>K}\sigma(s,S_s)S_s dW^S_s$ is a true martingale, we get
$$ \partial_tC_m(t,K)= \mathbb{E}^d\big[ D^d_{0t}\mathbb{I}_{S_t>K}(r_t^d K- r_t^f S_t) \big] + \frac{1}{2}\sigma(t,K)^2K^2 \mathbb{E}^d\big[ D^d_{0t}\delta_K(S_t) \big]$$
where $C_m(t,K)= \mathbb{E}^d[\mathcal{P}_t]$ is the price of the call option in the model.\\
By using that $\partial_K^2C_m(t,K)= \mathbb{E}^d\big[ D^d_{0t}\delta_K(S_t) \big] $, we rewrite
\begin{align*}
\partial_tC_m(t,K) = &\ \mathbb{E}^d\big[ D^d_{0t}\mathbb{I}_{S_t>K}(r_t^d K- r_t^f S_t) \big] + \frac{1}{2}\sigma(t,K)^2K^2\partial_K^2C_m(t,K)\\
= &\ \mathbb{E}^d\big[ D^d_{0t}\mathbb{I}_{S_t>K}(r_t^d - r_d^0 )K \big] + \mathbb{E}^d\big[ D^d_{0t}\mathbb{I}_{S_t>K} r_d^0 K \big] \\
&-\mathbb{E}^d\big[ D^d_{0t}\mathbb{I}_{S_t>K}(r_t^f - r_f^0 )S_t \big] - \mathbb{E}^d\big[ D^d_{0t}\mathbb{I}_{S_t>K} r_f^0 S_t \big]\\
& + \frac{1}{2}\sigma(t,K)^2K^2\partial_K^2C_m(t,K)
\end{align*}

We now rewrite the term in $r^f$ under the foreign measure. From section 2.8 in \cite{clark}, one has that the Radon–Nikodym derivative of $\mathbb{P}^f$ with respect to $\mathbb{P}^d$ is
$$ \frac{d\mathbb{P}^f}{d\mathbb{P}^d}=\exp\bigg( \int_0^t\sigma(s,S_s) dW^S_s - \frac{1}{2}\int_0^t\sigma(s,S_s)^2 ds\bigg)$$
Hence:
\begin{align*}
& \mathbb{E}^d\big[ D^d_{0t}\mathbb{I}_{S_t>K}(r_t^f - r_f^0 )S_t \big] \\
& = \mathbb{E}^d\bigg[ D^d_{0t}\mathbb{I}_{S_t>K}(r_t^f - r_f^0 )S_0\exp\bigg( \int_0^t (r_s^d - r_s^f)ds + \int_0^t\sigma(s,S_s) dW^S_s - \frac{1}{2}\int_0^t\sigma(s,S_s)^2 ds \bigg) \bigg] \\
& =  S_0\mathbb{E}^d\bigg[\mathbb{I}_{S_t>K}(r_t^f - r_f^0 )\exp\bigg( -\int_0^t r_s^fds\bigg)\cdot\exp\bigg( \int_0^t\sigma(s,S_s) dW^S_s - \frac{1}{2}\int_0^t\sigma(s,S_s)^2 ds \bigg)  \bigg] \\
& =  S_0\mathbb{E}^d\bigg[D^f_{0t}\mathbb{I}_{S_t>K}(r_t^f - r_f^0 )\frac{d\mathbb{P}^f}{d\mathbb{P}^d}\bigg]\\
& =  S_0\mathbb{E}^f\big[D^f_{0t}\mathbb{I}_{S_t>K}(r_t^f - r_f^0 )\big]
\end{align*}
Together with the fact that $\partial_K C_m(t,K)= -\mathbb{E}^d\big[ D^d_{0t}\mathbb{I}_{S_t>K} \big] $, one has
\begin{align*}
\partial_tC_m(t,K) = & \ K\mathbb{E}^d\big[ D^d_{0t}\mathbb{I}_{S_t>K}(r_t^d - r_d^0 )\big] -  S_0\mathbb{E}^f\big[D^f_{0t}\mathbb{I}_{S_t>K}(r_t^f - r_f^0 )\big]\\
&- K \partial_K C_m(t,K)r_d^0 - r_f^0\mathbb{E}^d[D^d_{0t}\mathbb{I}_{S_t>K} S_t]+ \frac{1}{2}\sigma(t,K)^2K^2\partial_K^2C_m(t,K)\\
= & \ K\mathbb{E}^d\big[ D^d_{0t}\mathbb{I}_{S_t>K}(r_t^d - r_d^0 )\big] -  S_0\mathbb{E}^f\big[D^f_{0t}\mathbb{I}_{S_t>K}(r_t^f - r_f^0 )\big]\\
& - r_d^0 K \partial_K C_m(t,K) - r_f^0\mathbb{E}^d[D^d_{0t}(S_t-K)^+]- r_f ^0K\mathbb{E}^d[D^d_{0t}\mathbb{I}_{S_t>K}]\\
& + \frac{1}{2}\sigma(t,K)^2K^2\partial_K^2C_m(t,K)\\
= & \ K\mathbb{E}^d\big[ D^d_{0t}\mathbb{I}_{S_t>K}(r_t^d - r_d^0 )\big] -  S_0\mathbb{E}^f\big[D^f_{0t}\mathbb{I}_{S_t>K}(r_t^f - r_f^0 )\big]\\
& - (r_d^0 - r_f^0) K \partial_K C_m(t,K) - r_f^0 C_m(t,K)  \\
& + \frac{1}{2}\sigma(t,K)^2K^2\partial_K^2C_m(t,K)
\end{align*}
So the model is calibrated to the market smile if and only if 
\begin{align*}
\partial_tC(t,K) =  & \ K\mathbb{E}^d\big[ D^d_{0t}\mathbb{I}_{S_t>K}(r_t^d - r_d^0 )\big] -  S_0\mathbb{E}^f\big[D^f_{0t}\mathbb{I}_{S_t>K}(r_t^f - r_f^0 )\big]\\
& - (r_d^0 - r_f^0) K \partial_K C(t,K) - r_f^0 C(t,K)  \\
& + \frac{1}{2}K^2\sigma(t,K)^2\partial_K^2C(t,K)
\end{align*}
The result follows easily from the definition of $\sigma_{Dup}(t,K)$ in (\ref{eq: threeFactorDupireVolatility}). For sufficiency, fix the law of the existing
model and hence its rate-tail forcing term. The asserted identity makes the
market surface solve the same inhomogeneous linear forward equation as $C_m$,
with coefficient $\sigma^2$ and this same forcing. Their difference solves the
homogeneous equation with zero initial and boundary data; the assumed
uniqueness gives $C_m=C$.
\end{proof}

\begin{rem}
The equivalence of (\ref{eq: threeFactorCalibrationCondition}) and (\ref{eq: singleMeasure}) is the identity $D^d_{0t}S_t=S_0D^f_{0t}\,\frac{d\mathbb{Q}^f}{d\mathbb{Q}^d}$ used in the proof. In the case of absence of a foreign rate (i.e. $r^f_t\equiv 0$), we get the usual calibration condition in (\ref{eq:calibration_condition_spot}).
\end{rem}

The form (\ref{eq: singleMeasure}) is the one to use numerically: both rate corrections are expectations under the same measure. Restated under the domestic $T$-forward measure $\mathbb{Q}^{d^T}$ it reads
\begin{equation}\label{eq: threeFactorCalibrationConditionForward}
\sigma(t,K)^2 = \sigma_{Dup}(t,K)^2  - P^d_{0T}\frac{ \mathbb{E}^{\mathbb{Q}^{d^T}}\big[ (P^d_{tT})^{-1}\big(K(r_t^d-r^0_d)-S_t(r_t^f-r^0_f)\big)\mathbb{I}_{S_t>K} \big]}{ \frac{1}{2}K^2\partial^2_KC(t,K) }
\end{equation}

\subsection{The Particle Method}
From Lemma 5.2.2 in \cite{hybridModel}, the model (\ref{eq: threeFactorModel}) under the domestic $T$-forward measure $\mathbb{Q}^{d^T}$ reads:
\begin{align*}
 \frac{df_t}{f_t}   = &\ \sigma(t,S_t)dW^{S^T}_t - \sigma_{P_{tT}}^d dB_t^{d^T} + \sigma_{P_{tT}}^f dB_t^{f^T} \\
 dr_t^d  = &\  (\mu^d_t + \sigma_t^d\sigma_{P_{tT}}^d)dt + \sigma_t^d dB_t^{d^T} \numberthis \label{eq: threeFactorModelForward}\\
 dr_t^f  = & \ \big(\mu^f_t - \rho_{Sf}\sigma(t,S_t)\sigma_t^f +\rho_{df}\sigma_t^f\sigma_{P_{tT}}^d \big)dt + \sigma_t^f dB_t^{f^T}
\end{align*}
where $f_t:=S_t P^f_{tT}/P^d_{tT}$ is the forward and $\sigma_{P_{tT}}^d$, $\sigma_{P_{tT}}^f$ the volatilities of the domestic $P^d_{tT}$ and foreign $P^f_{tT}$ bonds respectively.\\
A \textit{particle} is then described by the vector-valued process $(f_t^i,r_t^{d,(i)},r_t^{f,(i)})$, for $1\leq i \leq N$, where $N$ is the number of simulated particles. Using the representation (\ref{eq: threeFactorCalibrationConditionForward}) for the local volatility, define:
\begin{equation}\label{eq: discreteLocVol3FactorModel}
\sigma_N(t,K)^2  =  \sigma_{Dup}(t,K)^2  - P^d_{0T}\frac{\frac{1}{N}\sum_{i=1}^N (P_{tT}^{d,(i)})^{-1}\big(K(r_t^{d,(i)}-r^0_d)-S^i_t(r_t^{f,(i)}-r^0_f)\big)\mathbb{I}_{S_t^i>K} }{\frac{1}{2}K^2\partial^2_KC(t,K)}
\end{equation}
All particles are simulated under $\mathbb{Q}^{d^T}$: no second system of particles under the foreign forward measure is needed. The estimator (\ref{eq: discreteLocVol3FactorModel}) still has the two drawbacks discussed in Section \ref{subsec: malliavinGeneral}: it depends on $\partial_K^2C$, and it is a ratio of two small quantities far from the money. The Malliavin representation of Section \ref{subsec: malliavin3F} removes both from the rates correction. It does not remove them from the problem: $\sigma_{Dup}$ itself still requires $\partial_K^2C$, and the kernel regression still needs particles near the strike, so whether the wings behave better in a full calibration remains to be tested.

\subsection{Hull White Stochastic Interest Rates}\label{subsec: HW3F}
Assume both the domestic and foreign rates follow Hull-White dynamics as in (\ref{eq: hull_white}), with parameters $(\kappa_i,\sigma^i_r)$, $i\in \{d,f\}$. It is convenient to work with the centred states $x^i_t$ defined by
\begin{equation}\label{eq: centredStates}
r^i_t=r^0_i(t)+\frac{(\sigma^i_r)^2}{2\kappa_i^2}\big(1-e^{-\kappa_it}\big)^2+x^i_t,\qquad x^i_0=0,
\end{equation}
which under $\mathbb{Q}^d$ follow (linear SDEs, see for instance \cite{bruti})
\begin{align*}
dx^d_t&=-\kappa_dx^d_t\,dt+\sigma^d_r\,dB^d_t\\
dx^f_t&=\big(-\kappa_fx^f_t-\rho_{Sf}\sigma^f_r\,\sigma(t,S_t)\big)dt+\sigma^f_r\,dB^f_t
\end{align*}
In particular the centred rates
$$ c^i_t:=r^i_t-r^0_i(t)=x^i_t+\frac{(\sigma^i_r)^2}{2\kappa_i^2}\big(1-e^{-\kappa_it}\big)^2 $$
are known pathwise: the centring constant of the foreign rate is the market forward rate $r^0_f(t)=\mathbb{E}^{f^t}[r^f_t]$, which is \emph{not} $\mathbb{E}^{d^t}[r^f_t]$, and no expectation of $\sigma(s,S_s)$ has to be estimated.

The domestic bond price as a function of the domestic rate, and the foreign bond price as a function of the foreign rate, are given as in (\ref{eq: bondPrice_hullWhite}):
$$ P^i_{tT}= \frac{P_{0T}^{i,mkt}}{P_{0t}^{i,mkt}}\exp\bigg( -A^i(t,T)\,c^i_t-\frac{(\sigma^i_r)^2}{4\kappa_i^3}(e^{-\kappa_i T}-e^{-\kappa_i t})^2(e^{2\kappa_i t}-1) \bigg),\qquad A^i(t,T) = \frac{1-e^{ -\kappa_i(T-t) }}{\kappa_i}, $$
and the volatilities of the bonds are $\sigma_{P_{tT}}^i=-\sigma^i_r A^i(t,T)$. The Brownian motions of $\mathbb{Q}^d$ and $\mathbb{Q}^{d^T}$ are related by
\begin{align*}
dW^{S}_t = &\ dW^{S^T}_t - \rho_{Sd}\sigma^d_r A^d(t,T)dt \\
dB^{d}_t = &\ dB^{d^T}_t - \sigma^d_r A^d(t,T)dt  \numberthis \label{measureChange: domesticForwardToDomestic}\\
dB^{f}_t = &\ dB^{f^T}_t - \rho_{df}\sigma^d_r A^d(t,T)dt
\end{align*}
and the model (\ref{eq: threeFactorModelForward}) now reads, with $S_t=f_tP^d_{tT}/P^f_{tT}$:
\begin{align*}
 \frac{df_t}{f_t}   = &\ \sigma(t,S_t)dW^{S^T}_t + \sigma^d_r A^d(t,T) dB_t^{d^T} - \sigma^f_r A^f(t,T) dB_t^{f^T} \\
 dx_t^d  = &\  \big[-\kappa_dx^d_t - (\sigma^d_r)^2 A^d(t,T)\big]dt + \sigma^d_r dB_t^{d^T} \numberthis \label{eq: threeFactorModelForwardHullWhite}\\
 dx_t^f  = & \ \big[-\kappa_fx^f_t - \rho_{Sf}\sigma^f_r\sigma(t,S_t) - \rho_{df}\sigma^d_r\sigma^f_rA^d(t,T) \big]dt + \sigma^f_r dB_t^{f^T}
\end{align*}

\section{Malliavin representation with two stochastic rates}\label{subsec: malliavin3F}
Following \cite{guyonLabordere}, we give an alternative expression to the calibration condition (\ref{eq: threeFactorCalibrationCondition}), in order to achieve a more accurate and smooth estimation of the local volatility for strikes that are far from the money.

\subsection{The general formula}
Two elementary facts about the forward measures are needed.
\begin{lem}[Forward measures]\label{lem: forwardMeasures}
For every $t$, with $F_{0t}=S_0P^f_{0t}/P^d_{0t}$ the forward FX rate,
\begin{equation}\label{eq: RNforward}
\left.\frac{d\mathbb{Q}^{f^t}}{d\mathbb{Q}^{d^t}}\right|_{\mathcal{F}_t}=\frac{S_t}{F_{0t}} .
\end{equation}
Consequently, if $q^{d^t}$ and $q^{f^t}$ denote the densities of $S_t$ under the two measures,
\begin{itemize}
\item[(i)] $q^{f^t}(K)=\frac{K}{F_{0t}}\,q^{d^t}(K)$, and therefore $S_0P^f_{0t}\,q^{f^t}(K)=K\,P^d_{0t}\,q^{d^t}(K)=K\,\partial_K^2C(t,K)$;
\item[(ii)] $\mathbb{E}^{f^t}[H\,|\,S_t=K]=\mathbb{E}^{d^t}[H\,|\,S_t=K]$ for every integrable $\mathcal{F}_t$-measurable $H$.
\end{itemize}
\end{lem}
\begin{proof}
In domestic currency the numeraire of $\mathbb{Q}^{f^t}$ is $S_uP^f_{ut}$ and the numeraire of $\mathbb{Q}^{d^t}$ is $P^d_{ut}$; their ratio at $u=t$ is $S_t$, and $F_{0t}$ is its value at $u=0$. Since the density (\ref{eq: RNforward}) is a function of $S_t$ only, (i) is immediate and (ii) follows from Bayes' formula.
\end{proof}

\paragraph{Convention.} Write $(W^S,B^d,B^f)^\top=\mathcal{C}Z$ with $Z$ a three-dimensional standard Brownian motion and $\mathcal{C}\mathcal{C}^\top=R$. The Malliavin derivative with respect to the rate driver $B^d$ is $\MD^{B^d}_sF:=\langle \MD^Z_sF,\mathcal{C}^\top e_d\rangle$: it is the sensitivity of $F$ to a perturbation of the path of $B^d$ which moves $W^S$ and $B^f$ by the regression coefficients $\rho_{Sd}$ and $\rho_{df}$. With this convention $\mathbb{E}[F\int_0^th(s)dB^d_s]=\mathbb{E}[\int_0^th(s)\MD^{B^d}_sF\,ds]$. The same convention is used for the drivers of the other measures.

For this single-driver representation assume each rate is autonomous under its own forward measure and measurable with respect to its own rate-driver filtration; otherwise a full vector Clark--Ocone representation is required. By the Clark-Ocone formula, $r^d_t-r^0_d=\int_0^t\sigma^{d^t}_r(s)\,dB^{d^t}_s$ under $\mathbb{Q}^{d^t}$ and $r^f_t-r^0_f=\int_0^t\sigma^{f^t}_r(s)\,d\widetilde{B}^{f^t}_s$ under $\mathbb{Q}^{f^t}$, where $\sigma^{d^t}_r(s)=\mathbb{E}^{d^t}[\MD^{B^{d^t}}_sr^d_t|\mathcal{F}_s]$ and $\sigma^{f^t}_r(s)=\mathbb{E}^{f^t}[\MD^{\widetilde{B}^{f^t}}_sr^f_t|\mathcal{F}_s]$. The derivation of section 11.7.6.1 in \cite{guyonLabordere}, applied to each term of (\ref{eq: threeFactorCalibrationCondition}) under its own forward measure and combined with Lemma \ref{lem: forwardMeasures} (i), gives
\begin{align*}
\sigma(t,K)^2 = \sigma_{Dup}(t,K)^2 & - \frac{2}{K}\int _0^t \mathbb{E}^{d^t}\big[\sigma_r^{d^t}(s)\, \MD _s^{B^{d^t}}S_t  \big|S_t = K \big]ds \\
& +  \frac{2}{K}\int _0^t \mathbb{E}^{f^t}\big[\sigma_r^{f^t}(s)\, \MD _s^{\widetilde{B}^{f^t}}S_t  \big|S_t = K \big]ds
\numberthis \label{eq: threeFactorGeneralMalliavin}
\end{align*}
Both terms carry the factor $2/K$. In the foreign term,
$$S_0\,\mathbb{E}^f\big[D^f_{0t}\,\delta_K(S_t)\,\cdot\,\big]=S_0P^f_{0t}\,q^{f^t}(K)\,\mathbb{E}^{f^t}\big[\,\cdot\,\big|S_t=K\big],$$
and the density of $S_t$ under the foreign forward measure is $K\partial_K^2C/(S_0P^f_{0t})$, not $\partial^2_KC/P^f_{0t}$.

For Hull-White rates the Clark-Ocone integrands are deterministic,
\begin{equation}\label{eq: kernels}
\sigma_r^{d^t}(s)=\sigma^d_r\,e^{-\kappa_d(t-s)},\qquad \sigma_r^{f^t}(s)=\sigma^f_r\,e^{-\kappa_f(t-s)},
\end{equation}
because each rate is autonomous and Gaussian under its own forward measure. What remains is to evaluate the two integrated sensitivities of $S_t$.

\subsection{Why the foreign rate cannot be treated by a change of measure}\label{subsubsec: obstacle}
Under $\mathbb{Q}^{f}$ the foreign rate is autonomous, exactly as the domestic rate is under $\mathbb{Q}^d$, and it is tempting to compute its Malliavin derivative there and to transport the result to the domestic measure (a change-of-variable formula for the Malliavin derivative under a Girsanov shift is given in \cite{ewald}). This does not lead to a closed form, for two reasons.

First, the domestic term of (\ref{eq: threeFactorGeneralMalliavin}) requires the derivative with respect to $B^{d^t}$, not with respect to the foreign driver.
Second, and more fundamentally, by (\ref{measureChange: domesticToForeign}) the drivers of the two measures differ by $\rho_{Sf}\int_0^t\sigma(u,S_u)du$, which is a functional of the path: a perturbation of the noise moves $S$, hence the Girsanov kernel, hence $r^f$. One Brownian motion cannot be substituted for the other inside a Malliavin derivative. Indeed, differentiating the Hull-White foreign rate directly under $\mathbb{Q}^{d^T}$ gives
$$ d\,\MD^{B^{d^T}}_sr^f_t=\Big[-\kappa_f\,\MD^{B^{d^T}}_sr^f_t-\rho_{Sf}\,\sigma^f_r\,\partial_S\sigma(t,S_t)\,\MD^{B^{d^T}}_sS_t\Big]dt, $$
while $\MD^{B^{d^T}}_sS_t$ depends on $\MD^{B^{d^T}}_sr^f_t$ through the drift $r^d_t-r^f_t$ of the spot. The pair $(\MD_sS_t,\MD_sr^f_t)$ solves a coupled linear system with random coefficients; it has no scalar factorisation of the type $V_t/V_s$ that made the one-factor computation of Section \ref{subsec: malliavin1F} explicit. The symmetric statement holds under the foreign measure, where it is the domestic rate that carries a quanto drift $+\rho_{Sd}\sigma^d_r\sigma(t,S_t)$.

The way out is not to solve for the two-parameter family $\MD_sS_t$, $0\leq s\leq t$, at all: because the kernels (\ref{eq: kernels}) factorise as $e^{-\kappa t}\cdot e^{\kappa s}$, the $s$-integral in (\ref{eq: threeFactorGeneralMalliavin}) can be carried through the linear system.

\subsection{Integrated sensitivities}
Let $X_t=(S_t,r^d_t,r^f_t)^\top$. Under $\mathbb{Q}^d$, $dX_t=b(t,X_t)dt+\mathcal{A}(t,X_t)dZ_t$ with
$$ b(t,X)=\begin{pmatrix} S(r^d-r^f)\\ \kappa_d(\theta_d(t)-r^d)\\ \kappa_f(\theta_f(t)-r^f)-\rho_{Sf}\sigma^f_r\,\sigma(t,S)\end{pmatrix},\qquad
\mathcal{A}(t,X)=\begin{pmatrix} S\sigma(t,S)&0&0\\ 0&\sigma^d_r&0\\ 0&0&\sigma^f_r\end{pmatrix}\mathcal{C}. $$
Let $M_t=\nabla_Xb(t,X_t)$ and $G_{j,t}=\nabla_X\mathcal{A}_{\cdot j}(t,X_t)$ be the Jacobians of the drift and of the columns of $\mathcal{A}$, and define the integrated sensitivities
\begin{equation}\label{eq: integratedSensitivities}
Y^d_t:=\int_0^t\sigma^d_re^{-\kappa_d(t-s)}\,\MD^{B^{d^t}}_sX_t\,ds,\qquad
Y^f_t:=\int_0^t\sigma^f_re^{-\kappa_f(t-s)}\,\MD^{\widetilde{B}^{f^t}}_sX_t\,ds .
\end{equation}
\paragraph{Standing hypotheses for the Hull--White representation.}\label{hyp: representation}
Fix a finite horizon $T_*>0$, deterministic initial states as in
(\ref{eq: centredStates}), and a positive semidefinite correlation matrix $R$.
Choose any square root $\mathcal C\mathcal C^\top=R$; its inverse is not used.
The rate parameters are constant with $\kappa_i>0$. Assume:
\begin{itemize}
\item[(H1)] The coefficient $\sigma(t,S)$ is deterministic, continuous in $t$
and locally continuously differentiable in $S>0$. The state has a nonexplosive
strong solution with $S_t>0$. Its Cameron--Martin first variations exist under
the domestic and foreign risk-neutral measures and solve the differentiated
state equations. No global bound on $S\partial_S\sigma$ is imposed.
\item[(H2)] All numeraire-change stochastic exponentials are true martingales
on $[0,T_*]$. For each $t>0$, $S_t$ belongs to $\mathbb D^{1,2}$ under both
own $t$-forward measures. Its Malliavin derivative is square integrable on
$[0,t]\times\Omega$ under those measures, so integration against the bounded
Hull--White kernels and Malliavin duality are legitimate. The integrated
first variations and the drift and diffusion terms of their equations are
locally integrable to the orders needed for stochastic Fubini. In particular,
$\mathbb E^{d^t}|S_t\ell^d_t|+
\mathbb E^{f^t}|S_t\ell^f_t|<\infty$.
\item[(H3)] For the classical calibration statement, the discounted Tanaka
martingale is a true martingale and the rate-tail expectations are finite.
The model call surface is $C^{1,2}$ for $t>0,K>0$, and
$m(t,K)=\partial_{KK}C_m(t,K)>0$. The locally finite signed weighted measures
$\mathbb E^d[D^d_{0t}\ell^i_t\mathbf1_{S_t\in dK}]$ have continuous densities
$n_i(t,K)$. Conditional expectations at a strike use the continuous versions
$n_i(t,K)/m(t,K)$. These assumptions justify the indicator limit below.
\end{itemize}
The first two hypotheses suffice for the integrated system and smooth-payoff
duality. Hypothesis (H3) is needed only for the pointwise classical calibration
formula. Local differentiability alone is not asserted to imply (H2) or (H3).

\begin{prop}[Integrated sensitivities]\label{prop: integratedSensitivities}
Under the standing hypotheses, the processes $Y^d$ and $Y^f$ solve, with $Y^d_0=Y^f_0=0$,
\begin{align*}
dY_t^d &=\big[(M_t-\kappa_d I)Y_t^d+a_d(t,X_t)\big]dt +\sum_{j=1}^3G_{j,t}Y_t^d\,dZ_t^j \numberthis\label{eq: matrixSystem}\\
dY_t^f &=\big[(M_t+H_t-\kappa_f I)Y_t^f+a_f(t,X_t)\big]dt +\sum_{j=1}^3G_{j,t}Y_t^f\,dZ_t^j
\end{align*}
where $Z$ is the $\mathbb{Q}^d$-Brownian motion and
$$ a_d=\sigma^d_r\begin{pmatrix}\rho_{Sd}S\sigma\\ \sigma^d_r\\ \rho_{df}\sigma^f_r\end{pmatrix},\qquad
a_f=\sigma^f_r\begin{pmatrix}\rho_{Sf}S\sigma\\ \rho_{df}\sigma^d_r\\ \sigma^f_r\end{pmatrix},\qquad
H_t=\partial_S\sigma(t,S_t)\begin{pmatrix}S_t\sigma(t,S_t)\\ \rho_{Sd}\sigma^d_r\\ \rho_{Sf}\sigma^f_r\end{pmatrix}\begin{pmatrix}1&0&0\end{pmatrix}. $$
\end{prop}
\begin{proof}
For a deterministic $h$, $\int_0^th(s)\MD^{B}_sX_t\,ds$ is the derivative at $\varepsilon=0$ of the solution $X^\varepsilon_t$ obtained by shifting the driver $B$ by $\varepsilon\int_0^\cdot h(s)ds$ (and the correlated drivers accordingly): the Malliavin derivative in a Cameron-Martin direction is a directional derivative on the path space. Differentiating the SDE in $\varepsilon$ gives a linear SDE, the first variation equation with a source term.

\textit{Domestic.} For Hull-White rates $\mathbb{Q}^d$, $\mathbb{Q}^{d^t}$ and $\mathbb{Q}^{d^T}$ differ by deterministic drifts, (\ref{measureChange: domesticForwardToDomestic}), so shifting $B^{d^t}$ is the same as shifting the $\mathbb{Q}^d$-driver: $Z\to Z+\varepsilon\int h\,\mathcal{C}^\top e_d\,ds$. With $h(s)=\sigma^d_re^{\kappa_ds}$, the derivative $\mathcal{Y}_t=\partial_\varepsilon X^\varepsilon_t$ solves $d\mathcal{Y}_t=[M_t\mathcal{Y}_t+e^{\kappa_dt}a_d]dt+\sum_jG_{j,t}\mathcal{Y}_tdZ^j_t$, since $\mathcal{A}\,\mathcal{C}^\top e_d\,\sigma^d_r=a_d$. As $Y^d_t=e^{-\kappa_dt}\mathcal{Y}_t$, the first equation follows.

\textit{Foreign.} Again $\mathbb{Q}^f$ and $\mathbb{Q}^{f^t}$ differ by a deterministic drift, so the shift acts on the $\mathbb{Q}^f$-driver $\widetilde{Z}$, where by (\ref{measureChange: domesticToForeign}) $dZ_t=d\widetilde{Z}_t+\mathcal{C}^\top e_S\,\sigma(t,S_t)dt$. Under $\mathbb{Q}^f$ the drift of $X$ is $b+\sigma\,w$ with $w:=\mathcal{A}\,\mathcal{C}^\top e_S=(S\sigma,\rho_{Sd}\sigma^d_r,\rho_{Sf}\sigma^f_r)^\top$, so that the first variation equation has drift Jacobian $M_t+\nabla_X(\sigma w)$ and noise $\sum_jG_{j,t}\mathcal{Y}_td\widetilde{Z}^j_t=\sum_jG_{j,t}\mathcal{Y}_tdZ^j_t-\sigma\,\nabla_Xw\,\mathcal{Y}_tdt$. The two contributions of the Girsanov kernel combine into
$$ \nabla_X(\sigma w)-\sigma\nabla_Xw=w\,(\nabla_X\sigma)^\top=H_t . $$
The source term is $e^{\kappa_ft}a_f$ and $Y^f_t=e^{-\kappa_ft}\mathcal{Y}_t$.
\end{proof}
The matrix $H_t$ is the diffusion direction of the spot driver times the gradient of the Girsanov kernel between $\mathbb{Q}^d$ and $\mathbb{Q}^f$. It is exactly what a substitution of Brownian motions loses, and it is the reason why $Y^f$ does not satisfy the same variational system as $Y^d$. The systems (\ref{eq: matrixSystem}) have no index $s$ left, and do not depend on a terminal maturity.

\subsection{Reduction to four scalar processes}
Two components of (\ref{eq: matrixSystem}) are deterministic: $Y^{d,d}_t=v_d(t)$ and $Y^{f,f}_t=v_f(t)$, where
$$ v_i(t)=(\sigma^i_r)^2\,\frac{1-e^{-2\kappa_it}}{2\kappa_i},\qquad i\in\{d,f\}, $$
is the variance of the Hull-White state $x^i_t$ (for $Y^{f,f}$, note that the third row of $M_t+H_t$ is $(0,0,-\kappa_f)$: the quanto term cancels, as it must under $\mathbb{Q}^f$). For the remaining four, pass to logarithmic sensitivities of the spot,
$$ \ell^d_t:=\frac{Y^{d,S}_t}{S_t},\qquad \zeta^f_t:=Y^{d,f}_t,\qquad \ell^f_t:=\frac{Y^{f,S}_t}{S_t},\qquad \zeta^d_t:=Y^{f,d}_t , $$
and let $g(t,S)=S\,\partial_S\sigma(t,S)$. By It\^o's formula the terms $r^d-r^f$ and $\sigma\,dW^S$ of the spot cancel and (\ref{eq: matrixSystem}) becomes
\begin{align*}
d\ell^d_t &= \ell^d_t\Big[g(t,S_t)\big(dW^S_t-\sigma(t,S_t)dt\big)-\kappa_d\,dt\Big]
 +\Big[v_d(t)-\zeta^f_t+\rho_{Sd}\,\sigma^d_r\,\sigma(t,S_t)\Big]dt \\
d\zeta^f_t &= \Big[-(\kappa_d+\kappa_f)\zeta^f_t-\rho_{Sf}\,\sigma^f_r\,g(t,S_t)\,\ell^d_t
 +\rho_{df}\,\sigma^d_r\sigma^f_r\Big]dt \numberthis \label{eq: sensitivities}\\
d\ell^f_t &= \ell^f_t\Big[g(t,S_t)\,dW^S_t-\kappa_f\,dt\Big]
 +\Big[\zeta^d_t-v_f(t)+\rho_{Sf}\,\sigma^f_r\,\sigma(t,S_t)\Big]dt \\
d\zeta^d_t &= \Big[-(\kappa_d+\kappa_f)\zeta^d_t+\rho_{Sd}\,\sigma^d_r\,g(t,S_t)\,\ell^f_t
 +\rho_{df}\,\sigma^d_r\sigma^f_r\Big]dt
\end{align*}
with $\ell^d_0=\zeta^f_0=\ell^f_0=\zeta^d_0=0$ and $W^S$ the $\mathbb{Q}^d$-Brownian motion; under $\mathbb{Q}^{d^T}$ substitute $dW^S_t=dW^{S^T}_t-\rho_{Sd}\sigma^d_rA^d(t,T)dt$.

Two features of (\ref{eq: sensitivities}) are worth noting. The domestic pair has $g\,(dW^S-\sigma dt)$ where the foreign pair has $g\,dW^S$: this is the first component of $H_t$. The cross terms $-\rho_{Sf}\sigma^f_r\,g\,\ell^d$ and $+\rho_{Sd}\sigma^d_r\,g\,\ell^f$ are the feedback of the spot sensitivity into the rate that carries the quanto drift: this is the coupling of Section \ref{subsubsec: obstacle}. When the local volatility has no skew ($g=0$) everything is deterministic and $\zeta^f_t=\zeta^d_t$ is the covariance of $x^d_t$ and $x^f_t$.

\begin{lem}[Smooth-payoff duality]\label{lem: duality}
Under (H1)--(H2), for each bounded $C^1$ function $\psi$ with bounded derivative,
\begin{align*}
\mathbb E^d[D^d_{0t}c^d_t\psi(S_t)]
 &=\mathbb E^d[D^d_{0t}S_t\ell^d_t\psi'(S_t)],\\
\mathbb E^d[D^d_{0t}S_tc^f_t\psi(S_t)]
 &=\mathbb E^d[D^d_{0t}S_t^2\ell^f_t\psi'(S_t)],
\qquad c^i_t=r^i_t-r_i^0(t).
\end{align*}
No strike density or strict positive definiteness of $R$ is required here.
\end{lem}
\begin{proof}
Under the domestic $t$-forward measure,
$c_t^d=\int_0^t\sigma_r^d e^{-\kappa_d(t-s)}dB_s^{d^t}$.
The chain rule and duality against this deterministic square-integrable kernel
show that $\mathbb E^{d^t}[c_t^d\psi(S_t)]
=\mathbb E^{d^t}[S_t\ell_t^d\psi'(S_t)]$.
Multiplication by $P^d_{0t}$ gives the first identity. Repeat the argument
under the foreign $t$-forward measure and multiply by $S_0P^f_{0t}$.
The identity $S_0\mathbb E^f[D^f_{0t}H]=\mathbb E^d[D^d_{0t}S_tH]$
then gives the second identity. All changes of measure are applied to the
already-defined path functional $\ell^f$, not to an untransformed Malliavin operator.
\end{proof}

\begin{prop}[Malliavin representation, Hull-White rates]\label{prop: main}
Let both rates follow Hull-White dynamics and impose (H1)--(H3).
Let the market call surface $C$ have positive strike density and a positive
Dupire variance, with $C(0,K)=(S_0-K)^+$, $C(t,0)=S_0P^f_{0t}$ and
$C(t,K)\to0$ as $K\to\infty$. Assume uniqueness for its linear forward Dupire
PDE in a solution class containing both $C$ and the model call surface $C_m$,
with these initial and boundary conditions. Then the existing three-factor
model is calibrated to $C$ if and only if
\begin{equation}\label{eq: mainFormula}
\boxed{\;\sigma(t,K)^2=\sigma_{Dup}(t,K)^2-2\,\mathbb{E}^{d^t}\big[\ell^d_t-\ell^f_t\,\big|\,S_t=K\big]\;}
\end{equation}
where $\ell^d,\ell^f$ solve (\ref{eq: sensitivities}). Equivalently, for any $T\geq t$,
\begin{align*}
\sigma(t,K)^2&=\sigma_{Dup}(t,K)^2
-2\,\frac{\mathbb{E}^{\mathbb{Q}^{d^T}}\big[(P^d_{tT})^{-1}(\ell^d_t-\ell^f_t)\,\big|\,S_t=K\big]}{\mathbb{E}^{\mathbb{Q}^{d^T}}\big[(P^d_{tT})^{-1}\,\big|\,S_t=K\big]} \numberthis\label{eq: mainFormulaForward}\\
&=\sigma_{Dup}(t,K)^2
-2\,\frac{\mathbb{E}^{d}\big[D^d_{0t}(\ell^d_t-\ell^f_t)\,\big|\,S_t=K\big]}{\mathbb{E}^{d}\big[D^d_{0t}\,\big|\,S_t=K\big]} .
\end{align*}
\end{prop}
\begin{proof}
By (\ref{eq: kernels}) and (\ref{eq: integratedSensitivities}), the two integrals in (\ref{eq: threeFactorGeneralMalliavin}) are the conditional expectations of the first components $Y^{d,S}_t=S_t\ell^d_t$ and $Y^{f,S}_t=S_t\ell^f_t$, which equal $K\ell^d_t$ and $K\ell^f_t$ on $\{S_t=K\}$. The foreign conditional expectation is taken under $\mathbb{Q}^{f^t}$; by Lemma \ref{lem: forwardMeasures} (ii) it equals the one under $\mathbb{Q}^{d^t}$. The second line is the usual change of numeraire between $\mathbb{Q}^{d^t}$, $\mathbb{Q}^{d^T}$ and $\mathbb{Q}^d$.
For completeness, the calibration equivalence can be established without
identifying the model density with the market density in advance. Write
$m(t,K)=\mathbb E^d[D^d_{0t}\delta(S_t-K)]=\partial_{KK}C_m(t,K)$ and
$c^i_t=r^i_t-r_i^0(t)$. Apply Lemma \ref{lem: duality} to smooth approximations of
$\mathbf1_{(K,\infty)}$. Dominated convergence handles the left sides;
the continuous weighted densities in (H3) handle the approximate identities
on the right sides. This gives
\begin{align*}
\mathbb E^d[D^d_{0t}c^d_t\mathbb I_{S_t>K}]
 &=K m(t,K)\mathbb E^{d^t}[\ell^d_t\mid S_t=K],\\
\mathbb E^d[D^d_{0t}S_tc^f_t\mathbb I_{S_t>K}]
 &=K^2m(t,K)\mathbb E^{d^t}[\ell^f_t\mid S_t=K].
\end{align*}
Here the first identity uses the domestic forward measure; the second uses
foreign duality followed by Lemma \ref{lem: forwardMeasures}. It\^o--Tanaka
therefore yields, for the model surface,
\[
\partial_t C_m=-(r_d^0-r_f^0)K\partial_K C_m-r_f^0C_m
+\tfrac12K^2\left(\sigma^2(t,K)+2\mathbb E^{d^t}
[\ell^d_t-\ell^f_t\mid S_t=K]\right)\partial_{KK}C_m.
\]
If $C_m=C$, division by the positive market density proves necessity.
Conversely, substituting (\ref{eq: mainFormula}) makes this precisely the
market forward Dupire PDE. The initial conditions coincide. At $K=0$,
$C_m(t,0)=\mathbb E^d[D^d_{0t}S_t]=S_0P^f_{0t}$ by the true-martingale
numeraire change, and $C_m(t,K)\to0$ as $K\to\infty$ by integrability of
$D^d_{0t}S_t$. The assumed uniqueness therefore implies $C_m=C$. This argument proves neither existence
nor uniqueness of the nonlinear calibrated SDE.

\end{proof}
Compared with (\ref{eq: singleMeasure}), the correction no longer contains $\partial^2_KC(t,K)$ nor tail expectations, and it is a conditional expectation of quantities carried by the particles themselves.

\section{The Particle Method with Malliavin weights}\label{sec: algorithm3F}
A particle is described by $(f^i_t,x^{d,i}_t,x^{f,i}_t)$ together with the four scalars $(\ell^{d,i}_t,\zeta^{f,i}_t,\ell^{f,i}_t,\zeta^{d,i}_t)$, all simulated under $\mathbb{Q}^{d^T}$. Using (\ref{eq: mainFormulaForward}), define
\begin{equation}\label{eq: particleMalliavin3F}
\sigma_N(t,K)^2=\sigma_{Dup}(t,K)^2
-2\,\frac{\sum_{i=1}^N\Pi^i_t\,\big(\ell^{d,i}_t-\ell^{f,i}_t\big)\,\delta_{t,N}(S^i_t-K)}
{\sum_{i=1}^N\Pi^i_t\,\delta_{t,N}(S^i_t-K)},\qquad \Pi^i_t=\exp\big(A^d(t,T)\,x^{d,i}_t\big),
\end{equation}
where $\Pi^i_t$ is $(P^{d,i}_{tT})^{-1}$ up to a factor common to all particles.

\subsection{Time discretization of the sensitivities}
With the notation of Section \ref{subsec: discretisation3F}, $g_j=S_{t_j}\partial_S\sigma_N(t_j,S_{t_j})$, and the increment of the $\mathbb{Q}^d$-Brownian motion $\Delta W^{\mathbb{Q}^d}_j=\Delta W_j-\rho_{Sd}\sigma^d_r\,\Delta\,I^{1,d}_j$:
\begin{align*}
\ell^d_{t_{j+1}}&=\ell^d_{t_j}\big[1+g_j(\Delta W^{\mathbb{Q}^d}_j-\sigma_j\Delta)-\kappa_d\Delta\big]
+\big[v_d(t_j)-\zeta^f_{t_j}+\rho_{Sd}\sigma^d_r\sigma_j\big]\Delta\\
\zeta^f_{t_{j+1}}&=\zeta^f_{t_j}+\big[-(\kappa_d+\kappa_f)\zeta^f_{t_j}-\rho_{Sf}\sigma^f_r\,g_j\,\ell^d_{t_j}+\rho_{df}\sigma^d_r\sigma^f_r\big]\Delta \numberthis\label{eq: eulerSensitivities}\\
\ell^f_{t_{j+1}}&=\ell^f_{t_j}\big[1+g_j\,\Delta W^{\mathbb{Q}^d}_j-\kappa_f\Delta\big]
+\big[\zeta^d_{t_j}-v_f(t_j)+\rho_{Sf}\sigma^f_r\sigma_j\big]\Delta\\
\zeta^d_{t_{j+1}}&=\zeta^d_{t_j}+\big[-(\kappa_d+\kappa_f)\zeta^d_{t_j}+\rho_{Sd}\sigma^d_r\,g_j\,\ell^f_{t_j}+\rho_{df}\sigma^d_r\sigma^f_r\big]\Delta
\end{align*}

\subsection{The Algorithm}
Let $\{t_j\}$ be a discretization of $[0,T]$:
\begin{itemize}
\item[1.] Initialize $j=0$, set $\sigma_N(t,K)=\sigma_{Dup}(0,K)$ for all $t\in [t_0,t_1]$ and all sensitivities to zero
\item[2.] Simulate the $N$ particles  $\{(f_t^i, x_t^{d,i}, x^{f,i}_t) \}_{1\leq i\leq N}$ from  $t_{j}$ to  $t_{j+1}$ using  (\ref{eq: discretizationDomesticInterestRate}), (\ref{eq: discretizationForeignInterestRate}) and (\ref{eq: discretizationForward3F})
\item[3.] For each particle, update $(\ell^{d,i},\zeta^{f,i},\ell^{f,i},\zeta^{d,i})$ according to (\ref{eq: eulerSensitivities}), with the same Brownian increment
\item[4.] Compute $S^i_{t_{j+1}}=f^i_{t_{j+1}}P^{d,i}_{t_{j+1}T}/P^{f,i}_{t_{j+1}T}$ and sort the particles according to the spot value. For $K\in G_{K,t_{j+1}}$ compute $\sigma_N(t_{j+1},K)^2$ from (\ref{eq: particleMalliavin3F}), floored at a small positive value (Section \ref{subsec: caveats})
\begin{itemize}
\item Interpolate it and extrapolate it, with an interpolation that has a continuous first derivative in $K$, as $g_j$ requires it
\item Set $\sigma(t,K)\equiv \sigma_N(t_{j+1},K)$ for all $t\in [t_{j+1},t_{j+2}]$
\end{itemize}
\item[5.] Set $j:=j+1$, iterate steps 2,3 and 4 up to maturity T
\end{itemize}
The acceleration techniques of Section \ref{subsec: acceleration} apply unchanged.
\paragraph{Finite-particle interpretation.}
The sensitivities above differentiate a diffusion with deterministic coefficient
$\sigma$, including one determined self-consistently by a limiting law. At finite
$N$, the estimated coefficient depends on the entire population. The auxiliary
updates are a particle approximation to the limiting augmented system, not the
exact Malliavin derivatives of the finite interacting system: the latter would
also differentiate the empirical calibration function through other particles.
Convergence of this approximation, jointly with bandwidth and time-step limits,
is not established here.

\section{Numerical evidence}\label{sec: evidence}
\subsection{Consistency checks}\label{subsec: checks}

\paragraph{(a) No foreign rate.} If $\sigma^f_r=0$ then $\zeta^f=\zeta^d=\ell^f\equiv 0$ and the first line of (\ref{eq: sensitivities}) is (\ref{eq: oneProcess}): formula (\ref{eq: mainFormula}) is the one-factor Hull-White formula (\ref{eq: malliavinHullWhite}) of \cite{guyonLabordere}.

\paragraph{(b) Constant spot volatility.} If $\sigma(t,S)\equiv\sigma$ then $g=0$, the system (\ref{eq: sensitivities}) is deterministic and can be integrated:
$$ 2\big(\ell^d_t-\ell^f_t\big)=\partial_t\int_0^t\Big[
(\sigma^d_r A^d)^2+(\sigma^f_r A^f)^2+2\rho_{Sd}\sigma\sigma^d_r A^d-2\rho_{Sf}\sigma\sigma^f_r A^f
-2\rho_{df}\sigma^d_r\sigma^f_r A^dA^f\Big](s,t)\,ds . $$
The integral is the part of the total variance of the forward $f$ that is due to the rates in the Gaussian three-factor model, so the right hand side is exactly $\sigma_{Dup}(t)^2-\sigma^2$. With zero correlations this reads
$$ \sigma^2=\sigma_{Dup}(t)^2-\frac{(\sigma^d_r)^2}{\kappa_d^2}(1-e^{-\kappa_dt})^2-\frac{(\sigma^f_r)^2}{\kappa_f^2}(1-e^{-\kappa_ft})^2 . $$
This fixes all signs and the normalisation $2/K$ of both terms in (\ref{eq: threeFactorGeneralMalliavin}).

\paragraph{(c) Monte Carlo, with skew.} Integrating (\ref{eq: mainFormula}) twice in the strike gives two identities between unconditional expectations, which can be tested without any kernel estimate: with $c^i_T=r^i_T-r^0_i(T)$,
\begin{align*}
\mathrm{(I)}\quad&\mathbb{E}^d\big[D^d_{0T}\,c^d_T\,(S_T-K)^+\big]=\mathbb{E}^d\big[D^d_{0T}\,\mathbb{I}_{S_T>K}\,S_T\,\ell^d_T\big],\\
\mathrm{(II)}\quad&\mathbb{E}^d\big[D^d_{0T}\,S_T\,c^f_T\,(S_T-K)^+\big]=\mathbb{E}^d\big[D^d_{0T}\,\mathbb{I}_{S_T>K}\,S_T^2\,\ell^f_T\big].
\end{align*}
They were tested for a fixed local volatility with a pronounced skew,
$$\sigma(S)=0.12\,\big(1-0.6\tanh(2\ln S)\big),\qquad S_0=1,\qquad T=5,$$
with $\kappa_d=5\%$, $\sigma^d_r=1.2\%$, $\kappa_f=10\%$, $\sigma^f_r=1.5\%$, correlations
$\rho_{Sd}=0.35$, $\rho_{Sf}=-0.30$, $\rho_{df}=0.40$, flat curves at $3\%$ and
$1\%$, $2\cdot10^6$ antithetic paths and 200 Euler steps per year. In the table, $z$ is the difference between the two sides in standard errors. The column ``naive'' uses the system (\ref{eq: sensitivities}) without the cross terms and with $g(dW^S-\sigma dt)$ in $\ell^f$, i.e. without the corrections due to the measure change.

\begin{center}
\begin{tabular}{c|rrr|rrr|rr}
\toprule
 & \multicolumn{3}{c|}{(I) domestic, $\times10^{-3}$} & \multicolumn{5}{c}{(II) foreign, $\times10^{-3}$}\\
$K$ & direct & (\ref{eq: sensitivities}) & $z$ & direct & (\ref{eq: sensitivities}) & $z$ & naive & $z$\\
\midrule
0.70 & $2.4246$ & $2.4219$ & $+0.42$ & $-3.0956$ & $-3.0939$ & $-0.15$ & $-3.1798$ & $+6.96$\\
0.85 & $2.1273$ & $2.1244$ & $+0.52$ & $-2.8555$ & $-2.8537$ & $-0.18$ & $-2.9303$ & $+7.19$\\
1.00 & $1.7123$ & $1.7086$ & $+0.80$ & $-2.4536$ & $-2.4512$ & $-0.28$ & $-2.5123$ & $+6.98$\\
1.15 & $1.2363$ & $1.2361$ & $+0.04$ & $-1.9139$ & $-1.9138$ & $0.00$ & $-1.9562$ & $+6.66$\\
1.30 & $0.7868$ & $0.7876$ & $-0.35$ & $-1.3241$ & $-1.3245$ & $+0.08$ & $-1.3497$ & $+5.48$\\
1.50 & $0.3491$ & $0.3497$ & $-0.46$ & $-0.6578$ & $-0.6584$ & $+0.21$ & $-0.6684$ & $+3.65$\\
\bottomrule
\end{tabular}
\end{center}

\noindent Both identities hold within one standard error, while the naive foreign term is off by about $2.7\%$. For these parameters the cross term in the domestic pair is small (about $0.2\%$, not resolved by the test). In a stressed configuration ($T=10$, both rate volatilities doubled, skew coefficient $0.8$, 100 steps per year) both identities again hold within one standard error, the naive domestic term is off by about $1\%$ ($z$ between $3$ and $5$) and the naive foreign term by about $3\%$ ($z$ between $4$ and $8$).

The simulation scheme of Section \ref{subsec: discretisation3F} was checked separately: the closed forms of all step coefficients, including $K^{df}_j$, agree with numerical quadrature to $10^{-10}$; under $\mathbb{Q}^{d^T}$ the simulated particles reproduce $\mathbb{E}[f_t]=f_0$, $P^d_{0T}\mathbb{E}[(P^d_{tT})^{-1}]=P^d_{0t}$ and $P^d_{0T}\mathbb{E}[(P^d_{tT})^{-1}S_t]=S_0P^f_{0t}$ within Monte Carlo error; and (I)--(II) hold at an intermediate date ($t=4$, $T=6$) with $\mathbb{Q}^{d^T}$ particles and weights $(P^d_{tT})^{-1}$. Pooling five independent runs of $10^6$ paths at that date (two under $\mathbb{Q}^{d^T}$, three under $\mathbb{Q}^d$), the two sides differ at $K=0.7$ by $0.13\%$ for (I) and $0.09\%$ for (II), that is $0.8$ and $0.4$ standard errors; single runs deviate by up to $2.9$ standard errors at one strike, and the right hand sides computed under the two measures agree to about $0.1\%$.
The matrix form (\ref{eq: matrixSystem}), simulated directly, agrees with (\ref{eq: sensitivities}) path by path up to the Euler error. Scripts: \texttt{check\_sensitivities.py}, \texttt{check\_forward\_measure.py}, \texttt{check\_matrix\_form.py}.

\subsection{Calibration to a market surface}\label{subsec: cme}
The Monte Carlo evidence of Section \ref{subsec: checks} concerns the identities behind (\ref{eq: mainFormula}) for a \emph{given} smooth local volatility, and the simulation scheme. The full particle calibration has been run once, on CME EUR/USD option settlements of 12 November 2026: expiries from December 2026 to December 2027, the surface being extended beyond that at constant implied-volatility smile to five years, which is synthetic. The calibration is run under the domestic risk-neutral measure with discount weights, 200,000 particles and 100 time steps per year, and two sets of rate parameters which are illustrative and not calibrated to interest rate options. The options are then repriced with an independent simulation. Over the quoted expiries the repricing error attributable to the stochastic rates is at most 0.005 volatility points with the correction (\ref{eq: mainFormula}), against 0.19 and 0.47 for the two parameter sets with the Dupire local volatility alone. With the more volatile set, the calibrated variance reaches its floor within one standard deviation of the money from about three years, on the synthetic part of the surface, and the five-year at-the-money error is 0.54. The report and the scripts are in the repository \texttt{github.com/ceikit/sir-lsv-malliavin-particle}, file \texttt{calibration/RESULTS.md}.

\subsection{Synthetic comparison of three methods}\label{subsec: synthetic}
A second test is synthetic and uses options on the spot. A non-flat local volatility, whose level and skew vary in time, generates the target Dupire variance on a two-year horizon. The volatility is then recovered from that target by three methods: a three-dimensional forward equation for the discounted density with the conventional tail integrals, the five-field system of Section \ref{subsec: momentClosure}, and the particle method with the sensitivities (\ref{eq: sensitivities}). On their finest settings the three methods recover the generating volatility with a root-mean-square error between 0.002 and 0.003 volatility points, and the calibrated models reprice options against the generating model within 0.004. With 20,000 particles the particle method has an error of 0.004 after about three seconds; the two PDE solvers take 69 and 189 seconds on their finest grids. This is one set of parameters and one implementation of each method, and at this accuracy the error of the reference target is not negligible, so the methods cannot be ranked. Details are in \texttt{benchmarks/RESULTS.md}.

\subsection{What remains untested}
Several things remain untested. The CME options are options on futures and were treated as options on the forward to their expiry. The bias due to the time step and to the kernel bandwidth has not been quantified systematically. No market data beyond December 2027 was used, and no comparison has been made with other particle methods such as \cite{cozma,ogetbilHientzsch}.

\section{Existence and a deterministic formulation}\label{sec: theory}
\subsection{Weaker hypotheses and existence: precise scope}\label{subsec: existence}
The representation, existence for a prescribed coefficient, and existence of a
self-consistent calibrated coefficient are distinct assertions. The results below
separate them. The flat-smile construction is an exact calibration benchmark:
it permits a term structure but removes the spatial feedback through
$g=S\partial_S\sigma$. It is not a solution of the general-smile existence
problem. No novelty claim is made for these arguments.

\paragraph{What can be relaxed in the representation.}
Lemma \ref{lem: duality} requires neither a smooth positive strike density nor an
invertible correlation matrix. Local spatial differentiability and the stated
Malliavin integrability suffice; a global bound on $S\partial_S\sigma$ is not
needed. If the weighted measures have only locally integrable densities, the
indicator argument holds at their common Lebesgue points, so the strike identities
hold almost everywhere where $m(t,K)>0$. Continuous positive densities promote
these identities to all strikes. If no density is available, the smooth-payoff
identities remain valid, but we do not assert the classical pointwise Dupire
formula. We also do not assert the four-process equations for arbitrary
nondifferentiable local volatility without a separate weak-flow argument.

\begin{cor}[Deterministic time-dependent Gaussian rate parameters]
The representation and calibration equivalence extend to deterministic bounded
continuous $\kappa_i(t)\geq0$, $\sigma_i(t)\geq0$ and deterministic continuous
positive semidefinite correlation matrices $R(t)$, provided the corresponding
flow, integrability, density and PDE-uniqueness hypotheses hold. Write each
own-measure rate drift as $\alpha_i(t)-\kappa_i(t)r^i$ and fit the deterministic
$\alpha_i$ to its initial curve. In the four auxiliary equations replace all
parameters and correlations by their time-dependent values and replace $v_i$ by
\[
v_i(t)=\int_0^t\sigma_i(s)^2
          \exp\left(-2\int_s^t\kappa_i(u)du\right)ds,
\qquad \dot v_i=-2\kappa_i(t)v_i+\sigma_i(t)^2.
\]
The same calibration identity then holds.
\end{cor}
\begin{proof}
The own-forward Clark--Ocone kernel is
$\sigma_i(s)\exp(-\int_s^t\kappa_i(u)du)$, which factorizes into a function
of $s$ times a function of $t$. Use the Cameron--Martin direction
\[
h_i(s)=\sigma_i(s)\exp\left(\int_0^s\kappa_i(u)du\right).
\]
Rescale its first variation by $\exp(-\int_0^t\kappa_i(u)du)$.
Differentiating gives the same
matrix system with $-\kappa_i(t)Y^i$ and the time-dependent source. Correlation
square roots depend only on time, so their spatial derivatives add no terms.
Bond volatilities remain deterministic, and the foreign Girsanov correction
is still $w(\nabla\sigma)^\top$. The logarithmic reduction and duality proof
are unchanged. The rate's own component solves the displayed ODE for $v_i$.
\end{proof}
The covariance routines in this repository implement constant rate parameters;
this corollary is an analytical extension, not a claim of implemented support
for time-dependent parameters.

\begin{prop}[Local existence for the unregularized flat-smile problem]
\label{prop: gaussianExistence}
Let the market calls be Black FX calls with deterministic forward curves and
strike-independent total implied variance $V(t)$, where $V(0)=0$,
$V\in C^1([0,T_*])$ and $q(t):=V'(t)>0$. Fix constant Hull--White parameters
with $\kappa_d,\kappa_f>0$ and a valid correlation matrix. There exists $\tau>0$
and a deterministic positive spot volatility $a(t)$ on $[0,\tau)$ giving a
strong solution of the original, unregularized calibration equation
(\ref{eq: mainFormula}) and reproducing these calls.
The positive deterministic-volatility solution is unique on its maximal
interval of strict positivity. This is uniqueness within the Gaussian subclass,
not among all spot-dependent solutions.
\end{prop}
\begin{proof}
Put $B_i(t)=(1-e^{-\kappa_it})/\kappa_i$ and
\[
v_i(t)=\frac{\sigma_i^2(1-e^{-2\kappa_it})}{2\kappa_i},\qquad
z(t)=\rho_{df}\sigma_d\sigma_f
       \frac{1-e^{-(\kappa_d+\kappa_f)t}}{\kappa_d+\kappa_f},
\quad \sigma_i:=\sigma_r^i.
\]
Solve the two-dimensional nonautonomous ODE
\begin{align*}
\dot u_d&=-\kappa_du_d+v_d-z+\rho_{Sd}\sigma_d a,\\
\dot u_f&=-\kappa_fu_f+z-v_f+\rho_{Sf}\sigma_f a,\\
a(t)&=\sqrt{q(t)-2(u_d(t)-u_f(t))},\qquad u_d(0)=u_f(0)=0.
\end{align*}
On the open set $q(t)-2(u_d-u_f)>0$ the vector field is continuous in time
and locally Lipschitz in $(u_d,u_f)$. Since $q(0)>0$, the Picard--Lindelof
theorem gives a unique solution on a nonempty interval. The vector field has
at most linear growth in $|u_d|+|u_f|$, so there is no finite-time explosion
while inside this domain. Continuation holds on every compact time interval
where the radicand stays uniformly positive.

Set $\sigma(t,S)=a(t)$. The rate equations are linear with deterministic
coefficients and deterministic quanto drift; the log-spot equation is then an
explicit integral of these Gaussian rates plus a Gaussian stochastic integral.
This constructs a unique strong state process for this coefficient.
As $g=0$, the auxiliary system gives $\ell^d=u_d$, $\ell^f=u_f$ and
$\zeta^f=\zeta^d=z$, all deterministic. Consequently its conditional expectation
is $u_d-u_f$ under every forward measure and the original nonlinear identity
holds exactly, with no kernel, variance floor, or particle approximation.

To verify the market fit independently, under the domestic $t$-forward measure
$S_t$ is lognormal with mean $F_{0t}$ and log-variance
\begin{align*}
V_m(t)=\int_0^t\Big[&a(s)^2+
\sigma_d^2 A^d(s,t)^2+\sigma_f^2 A^f(s,t)^2\\
&+2a(s)\{\rho_{Sd}\sigma_d A^d(s,t)-\rho_{Sf}\sigma_f A^f(s,t)\}\\
&-2\rho_{df}\sigma_d\sigma_f A^d(s,t)A^f(s,t)\Big]ds.
\end{align*}
Differentiation, using $\partial_t A^i(s,t)=e^{-\kappa_i(t-s)}$
and the integrating-factor solutions of the two ODEs, gives
$V_m'(t)=a(t)^2+2(u_d-u_f)=q(t)$. Both variances vanish at zero, hence
$V_m=V$ and the Black call prices coincide. The ODE uniqueness proves the
asserted subclass uniqueness.
\end{proof}

\begin{cor}[Explicit global sufficient condition with zero spot--rate correlations]
If $\rho_{Sd}=\rho_{Sf}=0$, define
\[
Q(t)=\sigma_d^2 B_d(t)^2+\sigma_f^2 B_f(t)^2
       -2\rho_{df}\sigma_d\sigma_f B_d(t)B_f(t).
\]
Suppose $q(t)-Q(t)\geq\varepsilon>0$ throughout $[0,T_*]$.
Then the unregularized flat-smile problem has the positive solution
\[
a(t)=\sqrt{q(t)-Q(t)}
\]
on that entire horizon.
\end{cor}
\begin{proof}
The terms containing $a$ disappear from the auxiliary ODEs. Their solutions
satisfy $2(u_d-u_f)=Q(t)$, as follows either by integrating them or by
differentiating the deterministic rate contribution to $V_m$ above.
The positive gap gives continuation up to $T_*$.
\end{proof}
The condition is sufficient for this subclass. Failure of its positive gap
does not by itself exclude a spot-dependent calibrated solution.

\begin{prop}[A necessary log-contract compatibility condition]\label{prop: logbound}
Fix a maturity $t>0$. Suppose a calibrated solution has nonnegative spot
volatility, a square-integrable martingale part in its forward-log dynamics
on $[0,t]$, and finite
expected integrated forward variance. Define the target log-contract moment
\[
\mathcal V(t)=-2\mathbb E^{d^t}_{\rm target}
                 [\log(S_t/F_{0t})].
\]
Let $a_d(s)=\sigma_d A^d(s,t)$, $a_f(s)=\sigma_f A^f(s,t)$,
$R_t(s)=a_d(s)^2+a_f(s)^2-2\rho_{df}a_d(s)a_f(s)$, and
$b_t(s)=\rho_{Sd}a_d(s)-\rho_{Sf}a_f(s)$. Necessarily
\[
\mathcal V(t)\ \geq\ \int_0^t
\left[R_t(s)-\min\{b_t(s),0\}^2\right]ds.
\]
This bound applies to spot-dependent volatility as well as to the Gaussian
subclass, within the stated integrability class.
\end{prop}
\begin{proof}
For the forward $f_s=S_sP^f_{st}/P^d_{st}$, It\^o's formula under
$\mathbb Q^{d^t}$ and the integrability assumptions give
\[
-2\mathbb E^{d^t}\log(f_t/f_0)
=\mathbb E^{d^t}\int_0^t
   [\sigma(s,S_s)^2+2b_t(s)\sigma(s,S_s)+R_t(s)]ds.
\]
Since $f_t=S_t$ and $f_0=F_{0t}$, the left side equals the target moment.
Minimize the quadratic over a nonnegative volatility: its minimum is
$R_t(s)-\min\{b_t(s),0\}^2$. Integration gives the bound.
\end{proof}
For a flat smile, $\mathcal V(t)=V(t)$. In particular, with zero spot--rate
correlations any input with $V(t)<\int_0^tR_t(s)ds$ admits no calibrated
solution in this integrability class. An unconditional existence theorem for
arbitrary market surfaces and fixed rate parameters is therefore false.

\begin{prop}[Strong well-posedness of a specified regularized system]
\label{prop: regularizedExistence}
Use the domestic-measure augmented log-state
$U=(\log S,x^d,x^f,\log D^d,\ell^d,\zeta^f,\ell^f,\zeta^d)\in\mathbb R^8$.
Fix $h,\delta>0$ and a finite cutoff radius $L$. Choose a nonnegative kernel
$k\in C_b^3(\mathbb R)$ and write $k_h(x)=h^{-1}k(x/h)$.
Let $w_L(u)>0$ and $\lambda_L(u)$ be bounded $C^2$ functions with bounded
first and second derivatives, agreeing with $e^{u_4}$ and $u_5-u_7$ on
$[-L,L]^8$. Set
\[
M(t,y,\mu)=\frac{\int w_L(u)\lambda_L(u)k_h(u_1-y)\mu(du)}
 {\delta+\int w_L(u)k_h(u_1-y)\mu(du)},\qquad
s(t,y,\mu)=\sqrt{\chi\big(q_D(t,y)-2M(t,y,\mu)\big)}.
\]
Here $q_D(t,y)=\sigma_{Dup}^2(t,e^y)$ has uniformly bounded spatial
derivatives through order two and is continuous in time; $\chi\in C_b^3$
is a smooth variance limiter with values in a fixed interval
$[\underline a,\overline a]\subset(0,\infty)$.
Form the augmented SDE from the state and four auxiliary equations using
$s$ and $g=\partial_y s$. Evaluate its state arguments at a smooth bounded
coordinate cutoff equal to the identity on $[-L,L]^8$. Assume the
deterministic curve coefficients are bounded and continuous in time.
For every square-integrable initial state independent of the Brownian motion,
this explicitly regularized McKean--Vlasov SDE has a unique strong solution
on every finite horizon.
\end{prop}
\begin{proof}
The denominator defining $M$ is at least $\delta$. Bounded kernel derivatives
and bounded Lipschitz features imply that $M$, $\partial_yM$ and their
spatial Lipschitz constants are uniformly bounded. For any coupling of
$\mu,\nu$, integrating the feature differences gives bounds by a constant
times $W_1(\mu,\nu)\leq W_2(\mu,\nu)$, also for $\partial_yM$.
The quotient rule, positive variance floor and bounded derivatives of $\chi$
therefore make both $s$ and $g$ globally Lipschitz in $(y,\mu)$ for the
metric $|y-y'|+W_2(\mu,\nu)$. The cutoff bounds all remaining multiplicative
state factors in the augmented equations. Their drift $B$ and diffusion
$A$ consequently satisfy
\[
|B(t,u,\mu)-B(t,v,\nu)|+\|A(t,u,\mu)-A(t,v,\nu)\|
\leq C(|u-v|+W_2(\mu,\nu)),
\]
with bounded growth. The constants may depend on $h,\delta,L$.

For a deterministic continuous law flow, the frozen-law SDE has a unique
strong solution. Couple two such solutions with the same initial state and
Brownian motion. The usual It\^o/BDG estimate and Gronwall inequality yield
\[
\mathbb E\sup_{s\leq t}|U_s^\mu-U_s^\nu|^2
\leq C_T\int_0^t W_2(\mu_s,\nu_s)^2ds.
\]
Thus the frozen-law map is a contraction for sufficiently short time in the
supremum $W_2$ metric on continuous square-integrable law flows with fixed
initial law. Its fixed point, followed by successive time intervals, gives
strong existence on any finite horizon. For two solutions of the nonlinear
SDE the same coupling estimate, with
$W_2(\mathcal L(U_s),\mathcal L(V_s))^2\leq\mathbb E|U_s-V_s|^2$,
gives pathwise uniqueness by Gronwall; the fixed-law construction also gives
uniqueness in law.
\end{proof}
This theorem concerns a modified equation: kernel smoothing, denominator
regularization, bounded features, state cutoffs and a smooth variance limiter
are all part of its definition. It does not prove exact market calibration,
existence for the unregularized general-smile model, or convergence as
$h,\delta\downarrow0$ and $L\uparrow\infty$. Uniform estimates sufficient
for that limit have not been obtained here.

\paragraph{Route to the remaining general-smile problem.}
The single-measure tail formulation (\ref{eq: singleMeasure}) offers another
route to existence. Its denominator $K^2\partial_{KK}C/2$ is known from
market calls, and it contains no derivative of the unknown coefficient.
The obstacle moves to positivity, tail estimates and the discontinuous
indicator law functional. Proving an invariant positive-variance class and
compactness of its law map could provide a route to an unregularized existence
theorem. This is a research route, not a completed proof. Existing conditional
LSV existence results such as \cite{djete,mustapha} do not automatically supply
these estimates for the present stochastic-rate and discount-weighted system.

\subsection{Exact closure of discounted first moments}\label{subsec: momentClosure}
This section gives a deterministic formulation of the four-process
representation. The closure is exact for a prescribed, possibly nonconstant,
local volatility. A converse below identifies sufficiently regular solutions
of the nonlinear closed equations with calibrated stochastic models. It does
not assert existence of such nonlinear solutions.

\subsubsection{Coordinates and operators}
Use the domestic risk-neutral measure and set
\[
Z_t=(X_t,U_t,V_t)=(\log S_t,x_t^d,x_t^f),\qquad
s(t,x)=\sigma(t,e^x),\qquad g(t,x)=\partial_xs(t,x).
\]
Here $X$ denotes log spot, not the three-vector used in the manuscript's
variational construction. Write $\nu_i=\sigma_r^i$, $r_d=\eta_d(t)+u$,
$r_f=\eta_f(t)+v$, where the curve-fitting shifts are
\[
\eta_i(t)=r_i^0(t)+\tfrac12\nu_i^2 B_{\kappa_i}(t)^2,
\qquad B_k(t)=\int_0^t e^{-ka}da.
\]
In particular $B_0(t)=t$. Put
\[
\beta=\rho_{Sd}\nu_d,\quad \gamma=\rho_{Sf}\nu_f,\quad
c=\rho_{df}\nu_d\nu_f,\quad k=\kappa_d+\kappa_f,\quad
v_i(t)=\nu_i^2 B_{2\kappa_i}(t).
\]
The drift, covariance matrix and covariance with the spot Brownian motion are
\begin{equation}\label{eq: closureState}
b=\begin{pmatrix}r_d-r_f-\tfrac12s^2\\-\kappa_du\\-\kappa_fv-\gamma s\end{pmatrix},
\quad Q=\begin{pmatrix}s^2&\beta s&\gamma s\\\beta s&\nu_d^2&c\\\gamma s&c&\nu_f^2\end{pmatrix},
\quad h=\begin{pmatrix}s\\\beta\\\gamma\end{pmatrix}.
\end{equation}
Thus $d\langle Z,W^S\rangle_t=h(t,Z_t)dt$. Define the backward killed operator
and its distributional adjoint by
\begin{align}
\mathcal K_s\varphi&=b\cdot\nabla\varphi+
\tfrac12\sum_{a,b=1}^3Q_{ab}\partial_{ab}\varphi-r_d\varphi,
\label{eq: closureOperator}\\
\mathcal K_s^*w&=-\nabla\cdot(bw)+
\tfrac12\sum_{a,b=1}^3\partial_{ab}(Q_{ab}w)-r_dw.\nonumber
\end{align}
The off-diagonal second derivatives each occur twice inside the sum.

Order the auxiliary variables as $Y=(\ell^d,\zeta^f,\ell^f,\zeta^d)^\top$.
Their equations take the affine form
\begin{equation}\label{eq: closureAffine}
dY=(AY+f)dt+GY\,dW^S,
\end{equation}
where
\[
A=\begin{pmatrix}
-\kappa_d-sg&-1&0&0\\
-\gamma g&-k&0&0\\
0&0&-\kappa_f&1\\
0&0&\beta g&-k
\end{pmatrix},\qquad
f=\begin{pmatrix}v_d+\beta s\\c\\-v_f+\gamma s\\c\end{pmatrix},
\qquad G=\operatorname{diag}(g,0,g,0).
\]
The initial conditions are $Z_0=z_0=(\log S_0,0,0)$ and $Y_0=0$.

\subsubsection{The closed forward equations}
For $D_t=\exp(-\int_0^t r_d(a,Z_a)da)$ define a positive measure $p_t$
and four signed measures $m_{i,t}$ on $\mathbb R^3$ by
\begin{equation}\label{eq: closureMeasures}
\langle\varphi,p_t\rangle=\mathbb E^d[D_t\varphi(Z_t)],\qquad
\langle\varphi,m_{i,t}\rangle=\mathbb E^d[D_tY_{i,t}\varphi(Z_t)].
\end{equation}
These are discounted measures, not probability densities; $m_i$ need not be
positive. The same letters denote densities when they exist.

\begin{thm}[Exact discounted first-moment closure]\label{thm: momentClosure}
Suppose the state and affine auxiliary equations have nonexplosive solutions.
Assume $\mathbb E[D_t(1+|Y_t|)]<\infty$ and, for every compactly supported
smooth test function, integrability of the drift terms and the true-martingale
property of the local martingales in It\^o's formulas for
$D\varphi(Z)$ and $DY_i\varphi(Z)$. Then, in the distributional sense,
\begin{align}
\partial_t p&=\mathcal K_s^*p,\label{eq: closureFive}\\
\partial_t m_1&=\mathcal K_s^*m_1-(\kappa_d+sg)m_1-m_2
 +(v_d+\beta s)p-\nabla\cdot(ghm_1),\nonumber\\
\partial_t m_2&=\mathcal K_s^*m_2-km_2-\gamma gm_1+cp,\nonumber\\
\partial_t m_3&=\mathcal K_s^*m_3-\kappa_fm_3+m_4
 +(-v_f+\gamma s)p-\nabla\cdot(ghm_3),\nonumber\\
\partial_t m_4&=\mathcal K_s^*m_4-km_4+\beta gm_3+cp,\nonumber
\end{align}
with $p_0=\delta_{z_0}$ and $m_{i,0}=0$.
No second or higher moments of $Y$ occur. This assertion does not require
invertibility of the driver correlation matrix or existence of densities.
\end{thm}
\begin{proof}
The finite-variation differential $dD=-r_dDdt$ contributes the killing term.
The product rule for $DY_i\varphi(Z)$ gives the drift
\[
D\left\{Y_i\mathcal K_s\varphi+(AY+f)_i\varphi
            +(GY)_i h\cdot\nabla\varphi\right\}dt.
\]
The last summand is the quadratic covariation of $Y_i$ and $\varphi(Z)$.
Taking expectations gives
\[
\frac{d}{dt}\langle\varphi,m_i\rangle
=\langle\mathcal K_s\varphi,m_i\rangle+
\sum_j\langle A_{ij}\varphi,m_j\rangle+\langle f_i\varphi,p\rangle
+\mathbf1_{i\in\{1,3\}}\langle gh\cdot\nabla\varphi,m_i\rangle.
\]
Distributional integration by parts proves the assertion. Equivalently, the
joint generator preserves functions affine in $Y$. Their second derivatives
with respect to $Y$ vanish, so the quadratic auxiliary diffusion cannot produce
second auxiliary moments. All mixed state--auxiliary derivatives are retained
in the divergence term.
\end{proof}

\paragraph{Explicit shared-noise term.}
For either $m=m_1$ or $m=m_3$,
\[
-\nabla\cdot(ghm)=-\partial_x(sgm)-\beta\partial_u(gm)
                                      -\gamma\partial_v(gm).
\]
Omitting this term is incorrect even when the physical state diffusion is
otherwise implemented exactly. For fixed $s$ the five equations are linear:
one density equation and two pairs of coupled moment equations.

\subsubsection{Uniqueness for a prescribed nonflat coefficient}
The preceding theorem uses only integrability assumptions. The following
stronger, concrete assumptions give a convenient uniqueness class; they are
not claimed to be minimal.

\begin{prop}[Fixed-coefficient uniqueness]\label{prop: closureIdentification}
Fix a finite horizon. Assume the two rate volatilities are strictly positive,
the constant correlation matrix is positive definite, $\kappa_i\geq0$, and
the deterministic shifts are $C^1$.
Suppose $s$ is $C^{1,4}$, is bounded above and away from zero, and its spatial
derivatives through order four are bounded on the horizon. Then the stochastic
measures in (\ref{eq: closureMeasures}) are the unique weakly continuous
measure solution of (\ref{eq: closureFive}) with its specified initial data in
the class
\begin{equation}\label{eq: closureWeightClass}
\sup_{0\leq t\leq T}\int e^{\lambda|z|}
       \left(|p_t|+\sum_{i=1}^4|m_{i,t}|\right)(dz)<\infty
\quad\hbox{for every }\lambda>0.
\end{equation}
Here absolute values denote total variation. In particular the fixed-coefficient
system does not introduce additional solutions in this class. The volatility
may have genuine spot dependence.
\end{prop}
\begin{proof}
The state coefficients have linear growth and bounded spatial derivatives of
positive order needed below. The state SDE has a unique strong solution. The
rate states are Gaussian convolutions plus bounded drift convolutions; log
spot is an integral of these states plus a bounded-volatility martingale and
a bounded drift correction. Hence $Z$ has all exponential moments of $|Z|$,
uniformly on the finite horizon. The matrices $A,G$ and the source $f$ are
bounded. Standard It\^o moment estimates for their linear SDE give all finite
polynomial moments of $Y$. The domestic discount factor and its inverse have
all finite moments because the domestic integrated rate is Gaussian.
H\"older's inequality proves (\ref{eq: closureWeightClass}) for the stochastic
measures and the integrability required in Theorem \ref{thm: momentClosure}.

For uniqueness take arbitrary smooth compactly supported terminal functions
$\psi_0,\ldots,\psi_4$ at time $t$. Starting from $(Z_a,Y_a)=(z,y)$, define
\[
F(a,z,y)=\mathbb E_{a,z,y}\left[
 e^{-\int_a^t r_d(b,Z_b)db}
 \left(\psi_0(Z_t)+\sum_iY_{i,t}\psi_i(Z_t)\right)\right].
\]
The affine auxiliary equation implies
$F(a,z,y)=F_0(a,z)+\sum_i y_i F_i(a,z)$.
Differentiation of the stochastic flow up to order two is legitimate under the
bounded derivative assumptions; linear auxiliary moment estimates also apply
to these variations. It gives spatial derivatives of $F_j$ up to order two
with at most exponential growth in $|z|$. The Markov property and It\^o's
formula give the backward equations
\begin{align*}
\partial_a F_0+\mathcal K_s F_0+\sum_i f_iF_i&=0,\\
\partial_a F_j+\mathcal K_s F_j+\sum_i A_{ij}F_i
       +\mathbf1_{j\in\{1,3\}}gh\cdot\nabla F_j&=0.
\end{align*}
Here $f$ is the source in (\ref{eq: closureAffine}); the terminal conditions
are $F_j(t,\cdot)=\psi_j$.
Time regularity follows from these equations (or from the uniformly elliptic
backward equation). Pair these functions with the difference of two forward
solutions. Their time derivative is zero by duality. To use these noncompact
test functions, multiply by smooth cutoffs tending to one. All cutoff errors
vanish by (\ref{eq: closureWeightClass}), the exponential derivative bounds,
and the at-most-linear state drift and killing term. The initial difference
is zero, so the terminal pairing is zero for every choice of $\psi_j$.
This proves uniqueness of all five measures.
\end{proof}

\subsubsection{Compatibility identities and the spot marginal}
The next identities additionally use the fact that $Y$ consists of the
integrated Malliavin sensitivities, rather than arbitrary affine auxiliaries.
Impose the own-forward-measure duality and integrability assumptions of that
representation; the bounded smooth setting of Proposition
\ref{prop: closureIdentification} is a sufficient setting. Define
$c_d=r_d-r_d^0(t)$ and $c_f=r_f-r_f^0(t)$.

\begin{prop}[Full-state and marginal compatibility]\label{prop: closureCompatibility}
For the stochastic measures, the following are distributional identities:
\begin{align}
c_dp&=-\partial_xm_1-v_d\partial_up-\partial_vm_2,
\label{eq: closureConstraints}\\
c_fp&=-\partial_xm_3-m_3-\partial_um_4-v_f\partial_vp.\nonumber
\end{align}
Let bars denote integration over $(u,v)$, put $n_i=\overline{m_i}$,
$n=n_1-n_3$, and $J_i=\overline{c_i p}$. Then
\begin{equation}\label{eq: closureMarginalConstraints}
J_d=-\partial_xn_1,\qquad J_f=-\partial_xn_3-n_3.
\end{equation}
Consequently the discounted log-spot marginal obeys
\begin{equation}\label{eq: closureMarginal}
\partial_t\bar p=\mathcal B_t\bar p+
 \tfrac12(\partial_{xx}+\partial_x)(s^2\bar p+2n),
\qquad
\mathcal B_t w=-(r_d^0-r_f^0)\partial_xw-r_d^0 w.
\end{equation}
\end{prop}
\begin{proof}
For a smooth test $\varphi$ of all three coordinates, domestic own-forward
duality gives
\[
\mathbb E^d[D c_d\varphi(Z)]
=\mathbb E^d[D(\ell^d\partial_x\varphi+
                         v_d\partial_u\varphi+\zeta^f\partial_v\varphi)].
\]
Foreign own-forward duality gives the analogous identity with tangent vector
$(\ell^f,\zeta^d,v_f)$ and, on changing to the domestic measure, a factor
$S=e^x$ on both sides. Thus
\[
e^xc_fp=-\partial_x(e^xm_3)-\partial_u(e^xm_4)
                         -v_f\partial_v(e^xp).
\]
This proves (\ref{eq: closureConstraints}), including the essential $-m_3$
term. Rate integration, justified by cutoffs and the moment bounds, proves
(\ref{eq: closureMarginalConstraints}). Direct integration of the density
equation gives
\[
\partial_t\bar p=
-\partial_x\left[(r_d^0-r_f^0-\tfrac12s^2)\bar p+J_d-J_f\right]
+\tfrac12\partial_{xx}(s^2\bar p)-r_d^0\bar p-J_d.
\]
Substitution of the two marginal constraints yields
(\ref{eq: closureMarginal}).
\end{proof}

If the respective upper-boundary fluxes vanish, the same identities imply
\begin{equation}\label{eq: closureTailRecovery}
n_1(x)=\int_x^\infty J_d(y)dy,\qquad
n_3(x)=e^{-x}\int_x^\infty e^yJ_f(y)dy.
\end{equation}
Thus the correction can also be recovered from the physical discounted
density alone, using the conventional tail integrals. The moment fields give
an alternative evolution of those quantities; they do not add independent
terminal-law information. The five-field formulation is not a claim that
three-dimensional density-based calibration was previously unavailable.

\subsubsection{Nonlinear calibration and a known-denominator formulation}
Let $q(t,x)=\sigma_{Dup}(t,e^x)^2$ be the prescribed market Dupire variance.
With $\bar p>0$, the exact nonlinear feedback is
\begin{equation}\label{eq: closureFeedback}
s^2(t,x)=q(t,x)-2\frac{n(t,x)}{\bar p(t,x)},\qquad
n(t,x)=\int_{\mathbb R^2}(m_1-m_3)(t,x,u,v)du\,dv.
\end{equation}
There is no log-coordinate Jacobian in this ratio: it cancels between the
two marginals. Both are discounted under the domestic measure.
Equations (\ref{eq: closureFive}) and (\ref{eq: closureFeedback}), with
$g=\partial_xs$, are a closed nonlinear system of five fields on three spatial
coordinates. They are not five equations on log spot alone.

Let $p_*(t,x)$ denote the market discounted log-spot density. For a sufficiently
regular market call surface $C$,
\begin{equation}\label{eq: closureMarketDensity}
p_*(t,x)=e^x\partial_{KK}C(t,e^x),\quad
\partial_t p_* =\mathcal B_t p_*+
              \tfrac12(\partial_{xx}+\partial_x)(q p_*),\quad
p_*(0)=\delta_{\log S_0}.
\end{equation}
Assume it is positive for $t>0$ and that the stated one-dimensional forward
equations have uniqueness in the integrability class under consideration.
For smooth bounded uniformly positive coefficients and deterministic bounded
rates, this is the usual diffusion-law uniqueness class; exponential moments
also justify recovering call prices from the densities.

\begin{thm}[Calibration equivalence]
\label{thm: closureNonlinear}
Suppose the five fields solve (\ref{eq: closureFive}) in the class
(\ref{eq: closureWeightClass}), with densities for positive times and the
specified initial measure data. Suppose the induced deterministic coefficient $s$
satisfies Proposition \ref{prop: closureIdentification}, and the market data
satisfy (\ref{eq: closureMarketDensity}) and the uniqueness assumptions above.
Either of the feedback relations
\begin{align}
s^2&=q-2n/\bar p,\label{eq: closureTwoFeedbacks}\\
\text{or}\qquad s^2&=q-2n/p_*\label{eq: closureKnownDenominator}
\end{align}
for $t>0$ implies $\bar p=p_*$ and exact calibration of the associated
three-factor model. In this regularity class the two formulations are
equivalent, and their solutions are precisely the stochastic moment fields
of calibrated coefficients. This is an equivalence conditional on a solution,
not an existence or uniqueness assertion for the nonlinear system.
\end{thm}
\begin{proof}
Once the candidate $s$ is fixed, Proposition \ref{prop: closureIdentification}
identifies all five fields with the stochastic measures of the corresponding
strong state and auxiliary solution. Proposition \ref{prop: closureCompatibility}
therefore applies; it is not an additional freely imposed constraint.
With the first feedback, (\ref{eq: closureMarginal}) becomes exactly
(\ref{eq: closureMarketDensity}), so uniqueness gives $\bar p=p_*$.

For the second feedback use $2n=(q-s^2)p_*$. Subtract the market equation from
(\ref{eq: closureMarginal}) and put $w=\bar p-p_*$. Then
\begin{equation}\label{eq: closureDifference}
\partial_t w=\mathcal B_t w+
       \tfrac12(\partial_{xx}+\partial_x)(s^2w),\qquad w(0)=0.
\end{equation}
Uniqueness for this linear equation gives $w=0$. Hence the first feedback holds
as well. Finally
\[
C_m(t,K)=\int_{\mathbb R}(e^x-K)^+\bar p(t,x)dx=C(t,K).
\]
Conversely, a calibrated coefficient satisfying these assumptions has
$\bar p=p_*$ and satisfies the first feedback by the Malliavin calibration
identity. Its stochastic moments solve (\ref{eq: closureFive}), so they also
satisfy the second feedback.
\end{proof}

The known-denominator version removes division by an \emph{unknown} marginal
from this deterministic formulation. It does not remove small market densities
in the wings. Nor does replacing the marginal by a market density automatically
give an unbiased or consistent finite-particle estimator.

\paragraph{Analytical significance and remaining obstruction.}
The full discounted joint law of $(Z,Y)$ would require seven state coordinates.
For the calibration functional it can be replaced exactly by five fields on
the three physical coordinates. Discounting appears as killing, so an extra
discount-factor coordinate is unnecessary. For a prescribed coefficient,
Proposition \ref{prop: closureIdentification} gives a deterministic and
stochastic identification theorem with nonflat volatility.
For calibration, the known-denominator version is a useful formulation for
a fixed-point or perturbation argument. But it still contains a square root,
an unbounded rate domain, a time-zero Dirac law, and derivatives of the unknown
coefficient. In fact its divergence terms generally differentiate
$g=\partial_xs$ again. Positivity of the variance gap, control of small
densities and nonlinear estimates remain necessary. No contraction or
nonlinear parabolicity conclusion follows from the linear closure alone.

\paragraph{Checks and numerical scope.}
The companion script \texttt{checks/check\_moment\_closure.py} checks the
affine generator identity symbolically, checks both marginal PDE reductions,
and tests weak balances in a prescribed nonflat Monte Carlo model. These are
checks of the representation, not a three-dimensional PDE solver or a smile
recovery experiment. A solver should preserve discounted mass, monitor the
full-state compatibility residuals (\ref{eq: closureConstraints}), verify
the shared-noise fluxes, and compare recovered call prices independently.

\section{Limits and outlook}\label{sec: limits}
\subsection{Scope and caveats}\label{subsec: caveats}

\paragraph{Calibration identity versus existence of a calibrated model.} Propositions \ref{prop: calibration} and \ref{prop: main} are statements about a model that exists: \emph{if} (\ref{eq: threeFactorModel}) has a sufficiently regular solution, then it is calibrated if and only if its local volatility satisfies (\ref{eq: mainFormula}). Substituting (\ref{eq: mainFormula}) into (\ref{eq: threeFactorModel}) gives an SDE which is nonlinear in the sense of McKean, whose coefficient depends on the law of the solution through a conditional expectation given $S_t$. Neither existence nor uniqueness of a solution of this equation follows from the identity. The existence results above cover the stated Gaussian subclass and a regularized equation, not arbitrary unregularized smiles. Even for local-stochastic volatility models with deterministic rates, existence is known only in special cases: perturbative and local in time \cite{abergelTachet}, for regime-switching volatility under a condition on the range of the volatility factor \cite{jourdainZhou}, and for stationary solutions of a time-homogeneous version \cite{lackerShkolnikovZhang}. The more recent results of \cite{djete,mustapha} also concern different LSV structures and do not directly establish existence for the present general-smile stochastic-rates equation. The particle algorithm is therefore a numerical construction; its convergence as $N\to\infty$ is not proved in this paper.

\paragraph{Non-negativity of the local variance.}
The right hand side of (\ref{eq: mainFormula}) need not be positive.
The explicit Gaussian sufficient condition and the general necessary
log-contract bound in Section \ref{subsec: existence} have different scopes:
failure of a particular Gaussian construction is not a proof of nonexistence
among all spot-dependent models. In the numerical algorithm, report where
and by how much the variance floor binds. Such activation may reflect
estimation, interpolation or time-discretization error as well as incompatible
inputs; a negative trial iterate alone does not distinguish these cases.

\paragraph{Regularity.}
The hypotheses (H1)--(H3) distinguish local flow regularity and Malliavin
integrability from the density conditions needed for a classical strike-wise
formula. A global bound on $S\partial_S\sigma$ and a nonsingular correlation
matrix are not needed for smooth-payoff duality. In the algorithm, however,
noise in the estimated derivative $\partial_S\sigma_N$ enters the auxiliary
processes multiplicatively. Relaxing analytical assumptions does not remove
this numerical sensitivity.

\paragraph{Hull-White.} Proposition \ref{prop: main} uses two properties of the Hull-White model: the Clark-Ocone integrands (\ref{eq: kernels}) are deterministic and factorise, and the bond volatilities are deterministic, so that risk-neutral and forward measures differ by deterministic drifts.

\subsection{Outlook}
\paragraph{Existence.} In the system (\ref{eq: closureFive}), the equations of the two moments of $\ell^d$ and $\ell^f$ can be rewritten without second derivatives of the unknown variance $\sigma^2$ (a note in \texttt{paper/notes/} of the repository). Together with the known-denominator formulation of Theorem \ref{thm: closureNonlinear}, this is a possible starting point for a perturbative existence argument around the flat-smile solution of Proposition \ref{prop: gaussianExistence}. The estimates such an argument needs have not been obtained.

\paragraph{Gaussian multi-factor rates.} If a rate is a sum of Hull-White factors, the kernel (\ref{eq: kernels}) is a sum of exponentials, one per factor, and the matrix system (\ref{eq: matrixSystem}) extends with one Cameron-Martin direction per factor. The reduction to scalar processes has not been carried out for that case. Factors with different mean reversions cannot be aggregated into one process, and the cross sensitivities between each domestic and each foreign factor have to be tracked separately, so the number of auxiliary processes is not simply two per factor.

\paragraph{General one-factor short rates.} For rates $dr^i_t=\mu^i(t,r^i_t)dt+\sigma^i_r(t,r^i_t)dB^i_t$ the statement that replaces the factorisation $V_t/V_s$ is $\MD_sX_t=J_tJ_s^{-1}\mathcal{A}(s,X_s)\mathcal{C}^\top e$, with $J$ the $3\times3$ first variation matrix of $X=(S,r^d,r^f)$, which is not triangular because of the quanto drift. Two further difficulties appear: the Clark-Ocone integrand is stochastic and depends on $t$, and the forward measures differ from the risk-neutral ones by state-dependent drifts, which produce additional terms of the type of $H_t$.

\paragraph{Stochastic volatility.} With a stochastic volatility factor $a_t$ as in Section \ref{sec: oneRate}, the left hand side of (\ref{eq: mainFormula}) becomes $\sigma(t,K)^2\,\mathbb{E}^{d^t}[a_t^2|S_t=K]$ and the sensitivities acquire the tangent process of $a_t$ when $a_t$ is correlated with the rates. This is not worked out here.


\appendix
\section{One stochastic rate: time discretization and algorithm}\label{app: oneRate}
As anticipated before, at each iteration of the algorithm and for each particle, we need to simulate 3 processes
$f_t$, $r_t$, $V_t$, and compute 3 integrals $U_t$, $\Theta_t$, $\Lambda_t$.

\subsection{Processes: $f_t$ and $r_t$}
Let $\{t_j\}$ be a discretization of $[0,T]$ and consider a single particle.\\
In order to simulate $(f_{t_{j+1}}, r_{t_{j+1}} - r^0_{t_{j+1}})$ from $(f_{t_{j}}, r_{t_{j}} - r^0_{t_{j}})$, one can observe that, from (\ref{eq: hullWhiteSolution}):

\begin{align*}
r_{t_{j+1}}& - r^0_{t_{j+1}}  =  e^{-\kappa(t_{j+1}-t_j) }(r_{t_{j}} - r^0_{t_{j}})+\sigma_r\int_{t_j}^{t_{j+1}}e^{-\kappa(t_{j+1}-t) }dB_t^T \\
& - \frac{\sigma^2_r}{\kappa^2}e^{-\kappa t_{j+1} }\big( \sinh(\kappa t_{j+1})(1-e^{-\kappa(T-t_{j+1}) })-\sinh(\kappa t_{j})(1-e^{-\kappa(T-t_{j}) }) \big) \numberthis \label{eq: discretization_interestRate}
\end{align*}

and use the following approximation scheme for $f_t$:

\begin{align*}
f_{t_{j+1}}& = f_{t_{j}} \exp\bigg(  \sigma(t_j,P_{t_j T}f_{t_j})(W_{t_{j+1}}^T-W^T_{t_{j}})+\sigma_r\int_{t_j}^{t_{j+1}} \frac{1-e^{-\kappa(T-t)}}{\kappa}dB_t^T \\
& -\frac{1}{2}(t_{j+1}-t_{j})\big(\sigma(t_j,P_{t_j T}f_{t_j})^2+2\rho\sigma_r\sigma(t_j,P_{t_j T}f_{t_j})I^1_j + \sigma_r^2I_j^2\big)  \    \bigg) \numberthis \label{eq: discretization_forwardPrice}
\end{align*}
with 
\begin{align*}
& I^1_j = \frac{\kappa(t_{j+1}-t_{j})-e^{-\kappa T }(e^{\kappa t_{j+1} }-e^{\kappa t_{j} }) }{\kappa^2 (t_{j+1}-t_{j})} \\
& I^2_j = \frac{\kappa(t_{j+1}-t_{j})  -2e^{-\kappa T }(e^{\kappa t_{j+1} }-e^{\kappa t_{j} }) +\frac{1}{2}e^{-2\kappa T }(e^{2\kappa t_{j+1} }-e^{2\kappa t_{j} }) }{\kappa^3(t_{j+1}-t_{j})}
\end{align*}
Then the simulation of $(f_{t_{j+1}}, r_{t_{j+1}} - r^0_{t_{j+1}})$ boils down to simulate exactly the Gaussian vector
\[\bigg( W_{t_{j+1}}^T-W^T_{t_{j}},  \int_{t_j}^{t_{j+1}}e^{-\kappa(t_{j+1}-t) }dB_t^T  ,\int_{t_j}^{t_{j+1}} \frac{1-e^{-\kappa(T-t)}}{\kappa}dB_t^T \bigg)  \]
whose covariance matrix is $(t_{j+1}-t_{j})\Sigma_j$ where
\[\Sigma_j= \left( \begin{array}{ccc}
1 & \rho J_j^1 & \rho I_j^1 \\ 
\rho J_j^1 & J_j^2 & L_j \\ 
\rho I_j^1 & L_j & I^2_j \end{array} \right)\] 
with
\begin{align*}
& J_j^1 = \frac{1-e^{-\kappa(t_{j+1}-t_j) }}{\kappa(t_{j+1}-t_j)} \\
& J_j^2 = \frac{1-e^{-2\kappa(t_{j+1}-t_j) } }{2\kappa(t_{j+1}-t_j)} \\
& L_j = \frac{J_j^1}{\kappa}-e^{-\kappa(T + t_{j+1}) } \frac{ e^{2\kappa t_{j+1} }-e^{2\kappa t_{j} }  }{2\kappa^2(t_{j+1}-t_{j})}
\end{align*}
\subsection{Process: $V_t$}
For $V_t$ we shall start with a plain Euler discretization scheme. Refinements might follow.\\
The computation of the partial derivative $\partial_S\sigma$ in $V_t$ will be addressed in dependence on the chosen interpolation technique.
\subsection{Integrals: $U_t$, $\Theta_t$, $\Lambda_t$}
First note that $V_t$ is computed over the time grid $\{t_j\}$, hence it is a discrete vector of values.
Therefore we set $ V_t = (V_{t_{j+1}}+V_{t_{j}})/2 $ for all $t\in [t_j,t_{j+1}]$. Same reasoning applies for $\Theta_t$.\\
From (\ref{eq: malliavinWeights}) we have: 

\begin{align*}
\Theta_{t_{j+1}} & =  \int_0^{t_{j+1}}\frac{e^{-\kappa s}}{V_s}ds = \Theta_{t_{j}} + \int_{t_j}^{t_{j+1}}\frac{e^{-\kappa s}}{V_s}ds = 
\Theta_{t_{j}} + \frac{2}{(V_{t_{j+1}}+V_{t_{j}})}\int_{t_j}^{t_{j+1}}e^{-\kappa s}ds = \\
 & = \Theta_{t_{j}} - \frac{2}{(V_{t_{j+1}}+V_{t_{j}})}\frac{1}{\kappa}(e^{-\kappa t_{j+1}}-e^{-\kappa t_{j}} )\\ \\
\Lambda_{t_{j+1}} &= \int_0^{t_{j+1}}\sigma_re^{2\kappa s}\Theta_s ds = \Lambda_{t_{j}} + \int_{t_j}^{t_{j+1}}\sigma_re^{2\kappa s}\Theta_s ds = \Lambda_{t_{j}} + \frac{(\Theta_{t_{j+1}} +\Theta_{t_{j}} )}{2}\sigma_r\int_{t_j}^{t_{j+1}}e^{2\kappa s} =\\
& = \Lambda_{t_{j}} + \frac{(\Theta_{t_{j+1}} +\Theta_{t_{j}} )}{4\kappa}\sigma_r(e^{2\kappa t_{j+1}} - e^{2\kappa t_{j}})
\end{align*}
\\
A more careful analysis is needed for $U_t$ as it involves the integration of the local volatility function.
From (\ref{eq: malliavinWeights}) we have:
$$  U_{t_{j+1}} = \int_0^{t_{j+1}}e^{\kappa s}\frac{\sigma(s,S_s)}{V_s}ds =  U_{t_{j}} +  \int_{t_{j}}^{t_{j+1}}e^{\kappa s}\frac{\sigma(s,S_s)}{V_s}ds$$
From step 4 in the algorithm at section \ref{subsec: algorithm}, we have that $\sigma(t,K)\equiv \sigma_N(t_j,K)$ for all $t\in [t_j,t_{j+1}]$, hence we can rewrite the integral as: 
$$ U_{t_{j+1}} = U_{t_{j}} +  \frac{2}{(V_{t_{j+1}}+V_{t_{j}})}\int_{t_{j}}^{t_{j+1}}e^{\kappa s}\sigma(t_j,S_s)ds $$

\IfFileExists{volatility.jpg}{\centerline{\includegraphics[scale=0.5]{volatility.jpg}}}{\centerline{\fbox{\parbox{0.6\textwidth}{\centering\small [figure \texttt{volatility.jpg}: place the file next to this source]}}}}
\bigskip
Now we have two possibilities:
\begin{itemize}
\item approximate $\sigma(t_j,S_s)$ as $[ \sigma(t_j,S_{t_{j+1}}) + \sigma(t_j,S_{t_j}) ]/2$ or $\sigma(t_j,\frac{S_{t_{j+1}} +S_{t_{j}}}{2}) $ and analytically solve the remaining integral
\item apply numerical integration techniques to the integral 
\end{itemize}
The advantage of the first approach is that it is faster to compute and does not depend on the interpolation of the local volatility but it could imply non-negligible errors in the case $S_{t_j}$ is far from $S_{t_{j+1}}$.
The numerical integration might be slower to compute and one should pay attention at the interpolation algorithm used for the local volatility. However it provides a more accurate approximation of the integral, hence we may choose it above the former.

\subsubsection*{Trapezoidal rule for $U_t$}
If using a trapezoidal rule, $U_{t_{j+1}}$ reads:
$$ U_{t_{j+1}} = U_{t_{j}} +  \frac{(t_{j+1}-t_j)}{(V_{t_{j+1}}+V_{t_{j}})}\big[\sigma(t_j,S_{t_j})e^{\kappa t_j}+ \sigma(t_j,S_{t_{j+1}})e^{\kappa t_{j+1}} \big] $$

\subsubsection*{Cavalieri-Simpson rule for $U_t$}
If using a Cavalieri-Simpson rule, $U_{t_{j+1}}$ reads:
$$ U_{t_{j+1}} = U_{t_{j}} +  \frac{1}{3}\frac{(t_{j+1}-t_j)}{(V_{t_{j+1}}+V_{t_{j}})}\big[\sigma(t_j,S_{t_j})e^{\kappa t_j}+  4\sigma\big(t_j,\frac{S_{t_{j+1}} +S_{t_{j}}}{2}\big)e^{\kappa \frac{t_{j+1}+t_{j}}{2}}   +\sigma(t_j,S_{t_{j+1}})e^{\kappa t_{j+1}} \big] $$
\smallskip

\subsection{The Algorithm}\label{subsec: algorithm}
Let $\{t_j\}$ be a discretization of $[0,T]$:

\begin{itemize}
\item[1.] Initialize $j=0$ and set $\sigma_N(t,K)=\sigma_{Dup}(0,K)$ for all $t\in [t_0,t_1]$
\item[2.] Simulate the $N$ particles  $\{(f_t^i, r_t^i- r^0_t) \}_{1\leq i\leq N}$ from  $t_{j}$ to  $t_{j+1}$ using  (\ref{eq: discretization_interestRate}) and (\ref{eq: discretization_forwardPrice}) 
\item[3.] For each particle, simulate $V_{t_{j+1}}^i$ and compute the integrals $U_{t_{j+1}}^i$, $\Theta_{t_{j+1}}^i$, $\Lambda_{t_{j+1}}^i$ according to (\ref{eq: malliavinWeights})
\item[4.] Sort the particles according to the spot value. For $K\in G_{K,t_{j+1}}$, find the smallest index $i(K)$ and the largest $j(K)$ for which $\delta_{t,N}(S_t^i-K) > \eta$. Compute the local volatility $\sigma_N(t_{j+1},K)$ from (\ref{eq: discrete_localVolatility_forward_HullWhite}).
\begin{itemize}
\item Interpolate it and extrapolate it
\item Set $\sigma(t,K)\equiv \sigma_N(t_{j+1},K)$ for all $t\in [t_{j+1},t_{j+2}]$
\end{itemize}
\item[5.] Set $j:=j+1$, iterate steps 2,3 and 4 up to maturity T
\end{itemize}

\section{Time Discretization under $\mathbb{Q}^{d^T}$}\label{subsec: discretisation3F}
Let $\Delta=t_{j+1}-t_j$ and $\sigma_j=\sigma(t_j,S_{t_j})$. Freezing the local volatility over the step and integrating (\ref{eq: threeFactorModelForwardHullWhite}):
\begin{align*}
x^d_{t_{j+1}} ={}& e^{-\kappa_d\Delta}x^d_{t_j}+\sigma^d_r\,G^d_j-(\sigma^d_r)^2\Delta\,L^d_j \numberthis \label{eq: discretizationDomesticInterestRate}\\
x^f_{t_{j+1}} ={}& e^{-\kappa_f\Delta}x^f_{t_j}+\sigma^f_r\,G^f_j-\rho_{Sf}\sigma^f_r\,\sigma_j\,\Delta\,J^{1,f}_j-\rho_{df}\sigma^d_r\sigma^f_r\,\Delta\,L^{1,df}_j \numberthis \label{eq: discretizationForeignInterestRate}\\
f_{t_{j+1}} ={}&  f_{t_j}\exp\bigg( \sigma_j\,\Delta W_j + \sigma^d_r H^d_j - \sigma^f_r H^f_j
 - \frac{1}{2}\Delta\big[\sigma_j^2 + 2\rho_{Sd}\sigma^d_r\sigma_jI^{1,d}_j - 2\rho_{Sf}\sigma^f_r\sigma_jI^{1,f}_j   \\
& \qquad\qquad + (\sigma^d_r)^2 I^{2,d}_j + (\sigma^f_r)^2 I^{2,f}_j - 2\rho_{df} \sigma^d_r\sigma^f_rK^{df}_j\big]
\bigg) \numberthis \label{eq: discretizationForward3F}
\end{align*}
where
$$\Delta W_j=W^{S^T}_{t_{j+1}}-W^{S^T}_{t_{j}},\qquad G^i_j=\int_{t_j}^{t_{j+1}}e^{-\kappa_i(t_{j+1}-t) }dB_t^{i^T},\qquad H^i_j=\int_{t_j}^{t_{j+1}} A^i(t,T) dB_t^{i^T}. $$
The simulation of the three state variables boils down to simulate the Gaussian vector
$$\big(\Delta W_j,\,G^d_j,\,H^d_j,\,G^f_j,\,H^f_j\big)$$
whose covariance matrix is $\Delta\,\Sigma_j^{FX}$ where
\[\Sigma_j^{FX}= \left( \begin{array}{ccccc}
1 & \rho_{Sd} J_j^{1,d} & \rho_{Sd} I^{1,d}_j & \rho_{Sf} J_j^{1,f} & \rho_{Sf} I^{1,f}_j \medskip\\
\cdot & J_j^{2,d} & L_j^d & \rho_{df}J_j^{df} & \rho_{df}L^{2,df}_j \medskip\\
\cdot & \cdot & I^{2,d}_j & \rho_{df}L^{1,df}_j & \rho_{df}K^{df}_j \medskip\\
\cdot & \cdot & \cdot & J_j^{2,f} & L_j^f  \medskip\\
\cdot & \cdot & \cdot & \cdot & I^{2,f}_j
\end{array} \right)\]
with ($i\in \{d,f\}$)
\begin{align*}
& J_j^{1,i} = \frac{1-e^{-\kappa_i\Delta }}{\kappa_i\Delta}, \qquad
 J_j^{2,i} = \frac{1-e^{-2\kappa_i\Delta } }{2\kappa_i\Delta}, \qquad
 J_j^{df} = \frac{1 - e^{-(\kappa_f +\kappa_d)\Delta}}{(\kappa_f +\kappa_d)\Delta} \bigskip\\
& I^{1,i}_j = \frac{\kappa_i\Delta-e^{-\kappa_i T }(e^{\kappa_i t_{j+1} }-e^{\kappa_i t_{j} }) }{\kappa_i^2 \Delta} \\
& I^{2,i}_j = \frac{\kappa_i\Delta  -2e^{-\kappa_i T }(e^{\kappa_i t_{j+1} }-e^{\kappa_i t_{j} }) +\frac{1}{2}e^{-2\kappa_i T }(e^{2\kappa_i t_{j+1} }-e^{2\kappa_i t_{j} }) }{\kappa_i^3\Delta}\\
& L_j^i = \frac{J_j^{1,i}}{\kappa_i}-e^{-\kappa_i(T + t_{j+1}) } \frac{ e^{2\kappa_i t_{j+1} }-e^{2\kappa_i t_{j} }  }{2\kappa_i^2\Delta} \\
& L^{1,df}_j = \frac{J_j^{1,f}}{\kappa_d}-e^{- (\kappa_d T + \kappa_f t_{j+1}) } \frac{ e^{(\kappa_f +\kappa_d) t_{j+1} }-e^{(\kappa_f +\kappa_d) t_{j} }  }{\kappa_d(\kappa_f +\kappa_d)\Delta}\\
& L^{2,df}_j = \frac{J_j^{1,d}}{\kappa_f}-e^{- (\kappa_f T + \kappa_d t_{j+1}) } \frac{ e^{(\kappa_f +\kappa_d) t_{j+1} }-e^{(\kappa_f +\kappa_d) t_{j} }  }{\kappa_f(\kappa_f +\kappa_d)\Delta}\\
& K^{df}_j = \frac{1}{\kappa_d\kappa_f\Delta}\bigg[\Delta - \frac{e^{-\kappa_d T}(e^{\kappa_d t_{j+1}} - e^{\kappa_d t_j} )}{\kappa_d} - \frac{e^{-\kappa_f T}(e^{\kappa_f t_{j+1}} - e^{\kappa_f t_j} )}{\kappa_f}\\
& \qquad\qquad\qquad + \frac{e^{-(T-t_{j+1})(\kappa_f +\kappa_d)} - e^{-(T-t_{j})(\kappa_f +\kappa_d)}}{\kappa_f +\kappa_d}\bigg]
\end{align*}
The displayed closed forms identify the exact covariance, but direct
floating-point subtraction can be unstable for small mean reversion or small
steps. The accompanying implementation evaluates these averages by positive
Gauss--Legendre quadrature (at least 16 nodes), using
$A^i(u,T)=-\operatorname{expm1}(-\kappa_i(T-u))/\kappa_i$ and its continuous
limit $T-u$ at zero mean reversion. This numerical quadrature is not an additional
model time discretization; its accuracy is checked against independent dense trapezoidal
quadrature. It avoids cancellation in composite expressions. For sampling, QR factorization of the rectangular weighted kernel factor
is used directly, avoiding a Cholesky factorization of its rounded Gram matrix.
Extremely small steps can still require higher precision to resolve differences
below floating-point accuracy.

Each coefficient is the average over the step of a product of two of the kernels $1$, $e^{-\kappa_i(t_{j+1}-t)}$, $A^i(t,T)$; for instance
$$K^{df}_j=\frac{1}{\Delta}\int_{t_j}^{t_{j+1}}A^d(t,T)A^f(t,T)dt,\qquad L^{1,df}_j=\frac{1}{\Delta}\int_{t_j}^{t_{j+1}}A^d(t,T)e^{-\kappa_f(t_{j+1}-t)}dt.$$


\begin{thebibliography}{99}

\bibitem{guyonLabordere}
  Julien Guyon, Pierre Henry-Labord\`ere,
  \emph{Nonlinear Option Pricing},
  Chapman \& Hall/CRC, 2013.

\bibitem{deelstraRayee}
  Griselda Deelstra, Gr\'egory Ray\'ee,
  \emph{Local Volatility Pricing Models for Long-dated FX Derivatives},
  Applied Mathematical Finance 20(4), 380--402, 2013.

\bibitem{clark}
  Iain J. Clark,
  \emph{Foreign Exchange Option Pricing},
  2011.

\bibitem{nualart}
  David Nualart,
  \emph{The Malliavin Calculus and Related Topics},
  2005.

\bibitem{ewald}
  Christian-Oliver Ewald,
  \emph{A note on the Malliavin derivative operator under change of variable},
  2009.

\bibitem{hybridModel}
  Lech Aleksander Grzelak,
  \emph{Equity and Foreign Exchange Hybrid Models for Pricing Long-Maturity Financial Derivatives},
  2009.

\bibitem{bruti}
  Eckhard Platen, Nicola Bruti-Liberati,
  \emph{Numerical Solution of Stochastic Differential Equations with Jumps in Finance},
  2010.

\bibitem{ogetbilGanesanHientzsch}
  Orcan \"Ogetbil, Narayan Ganesan, Bernhard Hientzsch,
  \emph{Calibrating Local Volatility Models with Stochastic Drift and Diffusion},
  International Journal of Theoretical and Applied Finance 25(2), 2022. arXiv:2009.14764.

\bibitem{ogetbilHientzsch}
  Orcan \"Ogetbil, Bernhard Hientzsch,
  \emph{Extensions of Dupire Formula: Stochastic Interest Rates and Stochastic Local Volatility},
  arXiv:2005.05530, 2020.

\bibitem{cozma}
  Andrei Cozma, Matthieu Mariapragassam, Christoph Reisinger,
  \emph{Calibration of a Hybrid Local-Stochastic Volatility Stochastic Rates Model with a Control Variate Particle Method},
  SIAM Journal on Financial Mathematics 10(1), 181--213, 2019.

\bibitem{josephLoeperObloj}
  Benjamin Joseph, Gr\'egoire Loeper, Jan Ob{\l}\'oj,
  \emph{Calibration of Local Volatility Models with Stochastic Interest Rates using Optimal Transport},
  Finance and Stochastics 30, 397--439, 2026.

\bibitem{abergelTachet}
  Fr\'ed\'eric Abergel, R\'emi Tachet,
  \emph{A nonlinear partial integro-differential equation from mathematical finance},
  Discrete and Continuous Dynamical Systems - Series A 27(3), 907--917, 2010.

\bibitem{jourdainZhou}
  Benjamin Jourdain, Alexandre Zhou,
  \emph{Existence of a calibrated regime switching local volatility model},
  Mathematical Finance 30(2), 501--546, 2020.

\bibitem{lackerShkolnikovZhang}
  Daniel Lacker, Mykhaylo Shkolnikov, Jiacheng Zhang,
  \emph{Inverting the Markovian projection, with an application to local stochastic volatility models},
  Annals of Probability 48(5), 2189--2211, 2020.

\bibitem{djete}
Mao Fabrice Djete,
\emph{Non-regular McKean--Vlasov equations and calibration problem in local
stochastic volatility models}, arXiv:2208.09986v2, 2024.
\bibitem{mustapha}
Scander Mustapha,
\emph{Strong existence and uniqueness of a calibrated local stochastic
volatility model}, arXiv:2406.14074, 2024.
\bibitem{benhamouGruzRivoira}
E. Benhamou, A. Gruz and A. Rivoira,
\emph{Stochastic interest rates for local volatility hybrids models}, working paper, March 2008, SSRN 1107711.
\bibitem{blochNakashima}
D. Bloch and Y. Nakashima,
\emph{Multi-currency local volatility model}, working paper, June 2008 (revised August 2008), SSRN 1153337.
\end{thebibliography}


\end{document}
